% Options for packages loaded elsewhere
\PassOptionsToPackage{unicode}{hyperref}
\PassOptionsToPackage{hyphens}{url}
\PassOptionsToPackage{dvipsnames,svgnames,x11names}{xcolor}
%
\documentclass[
  11pt,
]{article}
\usepackage{amsmath,amssymb}
\usepackage{lmodern}
\usepackage{iftex}
\ifPDFTeX
  \usepackage[T1]{fontenc}
  \usepackage[utf8]{inputenc}
  \usepackage{textcomp} % provide euro and other symbols
\else % if luatex or xetex
  \usepackage{unicode-math}
  \defaultfontfeatures{Scale=MatchLowercase}
  \defaultfontfeatures[\rmfamily]{Ligatures=TeX,Scale=1}
\fi
% Use upquote if available, for straight quotes in verbatim environments
\IfFileExists{upquote.sty}{\usepackage{upquote}}{}
\IfFileExists{microtype.sty}{% use microtype if available
  \usepackage[]{microtype}
  \UseMicrotypeSet[protrusion]{basicmath} % disable protrusion for tt fonts
}{}
\makeatletter
\@ifundefined{KOMAClassName}{% if non-KOMA class
  \IfFileExists{parskip.sty}{%
    \usepackage{parskip}
  }{% else
    \setlength{\parindent}{0pt}
    \setlength{\parskip}{6pt plus 2pt minus 1pt}}
}{% if KOMA class
  \KOMAoptions{parskip=half}}
\makeatother
\usepackage{xcolor}
\IfFileExists{xurl.sty}{\usepackage{xurl}}{} % add URL line breaks if available
\IfFileExists{bookmark.sty}{\usepackage{bookmark}}{\usepackage{hyperref}}
\hypersetup{
  pdftitle={Applied Predictive Modeling --- Chapter 6 Solutions},
  pdfauthor={Randy Howk},
  colorlinks=true,
  linkcolor={blue},
  filecolor={Maroon},
  citecolor={Blue},
  urlcolor={blue},
  pdfcreator={LaTeX via pandoc}}
\urlstyle{same} % disable monospaced font for URLs
\usepackage[margin=0.5in]{geometry}
\usepackage{color}
\usepackage{fancyvrb}
\newcommand{\VerbBar}{|}
\newcommand{\VERB}{\Verb[commandchars=\\\{\}]}
\DefineVerbatimEnvironment{Highlighting}{Verbatim}{commandchars=\\\{\}}
% Add ',fontsize=\small' for more characters per line
\usepackage{framed}
\definecolor{shadecolor}{RGB}{248,248,248}
\newenvironment{Shaded}{\begin{snugshade}}{\end{snugshade}}
\newcommand{\AlertTok}[1]{\textcolor[rgb]{0.94,0.16,0.16}{#1}}
\newcommand{\AnnotationTok}[1]{\textcolor[rgb]{0.56,0.35,0.01}{\textbf{\textit{#1}}}}
\newcommand{\AttributeTok}[1]{\textcolor[rgb]{0.77,0.63,0.00}{#1}}
\newcommand{\BaseNTok}[1]{\textcolor[rgb]{0.00,0.00,0.81}{#1}}
\newcommand{\BuiltInTok}[1]{#1}
\newcommand{\CharTok}[1]{\textcolor[rgb]{0.31,0.60,0.02}{#1}}
\newcommand{\CommentTok}[1]{\textcolor[rgb]{0.56,0.35,0.01}{\textit{#1}}}
\newcommand{\CommentVarTok}[1]{\textcolor[rgb]{0.56,0.35,0.01}{\textbf{\textit{#1}}}}
\newcommand{\ConstantTok}[1]{\textcolor[rgb]{0.00,0.00,0.00}{#1}}
\newcommand{\ControlFlowTok}[1]{\textcolor[rgb]{0.13,0.29,0.53}{\textbf{#1}}}
\newcommand{\DataTypeTok}[1]{\textcolor[rgb]{0.13,0.29,0.53}{#1}}
\newcommand{\DecValTok}[1]{\textcolor[rgb]{0.00,0.00,0.81}{#1}}
\newcommand{\DocumentationTok}[1]{\textcolor[rgb]{0.56,0.35,0.01}{\textbf{\textit{#1}}}}
\newcommand{\ErrorTok}[1]{\textcolor[rgb]{0.64,0.00,0.00}{\textbf{#1}}}
\newcommand{\ExtensionTok}[1]{#1}
\newcommand{\FloatTok}[1]{\textcolor[rgb]{0.00,0.00,0.81}{#1}}
\newcommand{\FunctionTok}[1]{\textcolor[rgb]{0.00,0.00,0.00}{#1}}
\newcommand{\ImportTok}[1]{#1}
\newcommand{\InformationTok}[1]{\textcolor[rgb]{0.56,0.35,0.01}{\textbf{\textit{#1}}}}
\newcommand{\KeywordTok}[1]{\textcolor[rgb]{0.13,0.29,0.53}{\textbf{#1}}}
\newcommand{\NormalTok}[1]{#1}
\newcommand{\OperatorTok}[1]{\textcolor[rgb]{0.81,0.36,0.00}{\textbf{#1}}}
\newcommand{\OtherTok}[1]{\textcolor[rgb]{0.56,0.35,0.01}{#1}}
\newcommand{\PreprocessorTok}[1]{\textcolor[rgb]{0.56,0.35,0.01}{\textit{#1}}}
\newcommand{\RegionMarkerTok}[1]{#1}
\newcommand{\SpecialCharTok}[1]{\textcolor[rgb]{0.00,0.00,0.00}{#1}}
\newcommand{\SpecialStringTok}[1]{\textcolor[rgb]{0.31,0.60,0.02}{#1}}
\newcommand{\StringTok}[1]{\textcolor[rgb]{0.31,0.60,0.02}{#1}}
\newcommand{\VariableTok}[1]{\textcolor[rgb]{0.00,0.00,0.00}{#1}}
\newcommand{\VerbatimStringTok}[1]{\textcolor[rgb]{0.31,0.60,0.02}{#1}}
\newcommand{\WarningTok}[1]{\textcolor[rgb]{0.56,0.35,0.01}{\textbf{\textit{#1}}}}
\usepackage{graphicx}
\makeatletter
\def\maxwidth{\ifdim\Gin@nat@width>\linewidth\linewidth\else\Gin@nat@width\fi}
\def\maxheight{\ifdim\Gin@nat@height>\textheight\textheight\else\Gin@nat@height\fi}
\makeatother
% Scale images if necessary, so that they will not overflow the page
% margins by default, and it is still possible to overwrite the defaults
% using explicit options in \includegraphics[width, height, ...]{}
\setkeys{Gin}{width=\maxwidth,height=\maxheight,keepaspectratio}
% Set default figure placement to htbp
\makeatletter
\def\fps@figure{htbp}
\makeatother
\setlength{\emergencystretch}{3em} % prevent overfull lines
\providecommand{\tightlist}{%
  \setlength{\itemsep}{0pt}\setlength{\parskip}{0pt}}
\setcounter{secnumdepth}{5}
\usepackage{xcolor}
\usepackage{fvextra}
\DefineVerbatimEnvironment{Highlighting}{Verbatim}{breaklines,breakanywhere,commandchars=\\\{\}}
\usepackage{graphicx}
\usepackage{float}
\usepackage{microtype}
\usepackage{setspace}
\setstretch{1.05}
\setlength{\emergencystretch}{1em}
\makeatletter
\def\maxwidth{\ifdim\Gin@nat@width>\linewidth\linewidth\else\Gin@nat@width\fi}
\def\maxheight{\ifdim\Gin@nat@height>0.85\textheight0.85\textheight\else\Gin@nat@height\fi}
\makeatother
\setkeys{Gin}{width=\maxwidth,height=\maxheight,keepaspectratio}
\floatplacement{figure}{H}
\ifLuaTeX
  \usepackage{selnolig}  % disable illegal ligatures
\fi

\title{Applied Predictive Modeling --- Chapter 6 Solutions}
\author{Randy Howk}
\date{2026-04-01}

\begin{document}
\maketitle

{
\hypersetup{linkcolor=}
\setcounter{tocdepth}{2}
\tableofcontents
}
\hypertarget{how-to-use-this-file}{%
\section{How to use this file}\label{how-to-use-this-file}}

This R Markdown file contains \textbf{worked solutions} for all Chapter
6 exercises from \emph{Applied Predictive Modeling}: 6.1 (Tecator fat
prediction), 6.2 (permeability), and 6.3 (chemical manufacturing
process). The chapter covers ordinary linear regression, partial least
squares (PLS), and penalized linear models such as ridge, lasso, and
elastic net. The exercises appear at the end of Chapter 6 and ask the
reader to use those model families on three regression problems.
fileciteturn0file0

The code is written to be executable in a standard R setup with the
listed packages installed. Where the chapter asks for judgments such as
``which model would you choose?'', the text below gives a direct answer
and the code computes the relevant summaries.

\hypertarget{chapter-6-in-one-paragraph}{%
\section{Chapter 6 in one paragraph}\label{chapter-6-in-one-paragraph}}

Chapter 6 studies models that are linear in the parameters, emphasizing
three recurring ideas: (1) ordinary least squares is simple and
interpretable but can become unstable with correlated predictors or when
the number of predictors is large relative to sample size; (2) PLS is
often preferable to PCR because it uses the response while building
latent components; and (3) penalized methods such as ridge and lasso
improve prediction by shrinking coefficients and controlling variance.
The chapter also stresses preprocessing, resampling, and out-of-sample
evaluation rather than relying only on training-set fit.
fileciteturn3file1L8-L21 fileciteturn3file8L1-L21

\hypertarget{helper-functions}{%
\section{Helper functions}\label{helper-functions}}

\begin{Shaded}
\begin{Highlighting}[]
\NormalTok{rmse\_r2 }\OtherTok{\textless{}{-}} \ControlFlowTok{function}\NormalTok{(obs, pred) \{}
\NormalTok{  out }\OtherTok{\textless{}{-}} \FunctionTok{data.frame}\NormalTok{(}\AttributeTok{obs =}\NormalTok{ obs, }\AttributeTok{pred =}\NormalTok{ pred)}
\NormalTok{  caret}\SpecialCharTok{::}\FunctionTok{defaultSummary}\NormalTok{(out)}
\NormalTok{\}}

\NormalTok{summarize\_resamples }\OtherTok{\textless{}{-}} \ControlFlowTok{function}\NormalTok{(...) \{}
\NormalTok{  rs }\OtherTok{\textless{}{-}} \FunctionTok{resamples}\NormalTok{(}\FunctionTok{list}\NormalTok{(...))}
  \FunctionTok{list}\NormalTok{(}\AttributeTok{summary =} \FunctionTok{summary}\NormalTok{(rs), }\AttributeTok{diffs =} \FunctionTok{summary}\NormalTok{(}\FunctionTok{diff}\NormalTok{(rs)), }\AttributeTok{object =}\NormalTok{ rs)}
\NormalTok{\}}

\NormalTok{plot\_obs\_pred }\OtherTok{\textless{}{-}} \ControlFlowTok{function}\NormalTok{(obs, pred, }\AttributeTok{title =} \StringTok{"Observed vs Predicted"}\NormalTok{) \{}
  \FunctionTok{ggplot}\NormalTok{(}\FunctionTok{data.frame}\NormalTok{(}\AttributeTok{obs =}\NormalTok{ obs, }\AttributeTok{pred =}\NormalTok{ pred), }\FunctionTok{aes}\NormalTok{(}\AttributeTok{x =}\NormalTok{ obs, }\AttributeTok{y =}\NormalTok{ pred)) }\SpecialCharTok{+}
    \FunctionTok{geom\_point}\NormalTok{(}\AttributeTok{alpha =} \FloatTok{0.6}\NormalTok{, }\AttributeTok{shape =} \DecValTok{16}\NormalTok{) }\SpecialCharTok{+}
    \FunctionTok{geom\_abline}\NormalTok{(}\AttributeTok{slope =} \DecValTok{1}\NormalTok{, }\AttributeTok{intercept =} \DecValTok{0}\NormalTok{, }\AttributeTok{linetype =} \DecValTok{2}\NormalTok{, }\AttributeTok{linewidth =} \FloatTok{0.35}\NormalTok{, }\AttributeTok{color =} \StringTok{"gray45"}\NormalTok{) }\SpecialCharTok{+}
    \FunctionTok{labs}\NormalTok{(}\AttributeTok{title =}\NormalTok{ title, }\AttributeTok{x =} \StringTok{"Observed"}\NormalTok{, }\AttributeTok{y =} \StringTok{"Predicted"}\NormalTok{)}
\NormalTok{\}}

\NormalTok{plot\_tuning }\OtherTok{\textless{}{-}} \ControlFlowTok{function}\NormalTok{(model, x, }\AttributeTok{y =} \StringTok{"RMSE"}\NormalTok{, }\AttributeTok{color =} \ConstantTok{NULL}\NormalTok{, }\AttributeTok{title =} \ConstantTok{NULL}\NormalTok{) \{}
\NormalTok{  results }\OtherTok{\textless{}{-}} \FunctionTok{as\_tibble}\NormalTok{(model}\SpecialCharTok{$}\NormalTok{results)}

\NormalTok{  p }\OtherTok{\textless{}{-}} \FunctionTok{ggplot}\NormalTok{(results, }\FunctionTok{aes}\NormalTok{(}\AttributeTok{x =}\NormalTok{ .data[[x]], }\AttributeTok{y =}\NormalTok{ .data[[y]]))}

  \ControlFlowTok{if}\NormalTok{ (}\SpecialCharTok{!}\FunctionTok{is.null}\NormalTok{(color)) \{}
\NormalTok{    p }\OtherTok{\textless{}{-}}\NormalTok{ p }\SpecialCharTok{+}
      \FunctionTok{geom\_line}\NormalTok{(}\FunctionTok{aes}\NormalTok{(}\AttributeTok{color =} \FunctionTok{factor}\NormalTok{(.data[[color]]), }\AttributeTok{group =}\NormalTok{ .data[[color]])) }\SpecialCharTok{+}
      \FunctionTok{geom\_point}\NormalTok{(}\FunctionTok{aes}\NormalTok{(}\AttributeTok{color =} \FunctionTok{factor}\NormalTok{(.data[[color]]))) }\SpecialCharTok{+}
      \FunctionTok{labs}\NormalTok{(}\AttributeTok{color =}\NormalTok{ color)}
\NormalTok{  \} }\ControlFlowTok{else}\NormalTok{ \{}
\NormalTok{    p }\OtherTok{\textless{}{-}}\NormalTok{ p }\SpecialCharTok{+}
      \FunctionTok{geom\_line}\NormalTok{() }\SpecialCharTok{+}
      \FunctionTok{geom\_point}\NormalTok{()}
\NormalTok{  \}}

\NormalTok{  p }\SpecialCharTok{+} \FunctionTok{labs}\NormalTok{(}\AttributeTok{title =}\NormalTok{ title, }\AttributeTok{x =}\NormalTok{ x, }\AttributeTok{y =}\NormalTok{ y)}
\NormalTok{\}}

\NormalTok{vip\_pls }\OtherTok{\textless{}{-}} \ControlFlowTok{function}\NormalTok{(pls\_model) \{}
  \ControlFlowTok{if}\NormalTok{ (}\SpecialCharTok{!}\FunctionTok{inherits}\NormalTok{(pls\_model}\SpecialCharTok{$}\NormalTok{finalModel, }\StringTok{"mvr"}\NormalTok{)) }\FunctionTok{stop}\NormalTok{(}\StringTok{"Model must be a PLS model fit with caret::train(method = \textquotesingle{}pls\textquotesingle{})."}\NormalTok{)}
\NormalTok{  mod }\OtherTok{\textless{}{-}}\NormalTok{ pls\_model}\SpecialCharTok{$}\NormalTok{finalModel}
\NormalTok{  W }\OtherTok{\textless{}{-}}\NormalTok{ mod}\SpecialCharTok{$}\NormalTok{loading.weights}
\NormalTok{  T }\OtherTok{\textless{}{-}}\NormalTok{ mod}\SpecialCharTok{$}\NormalTok{scores}
\NormalTok{  Q }\OtherTok{\textless{}{-}}\NormalTok{ mod}\SpecialCharTok{$}\NormalTok{Yloadings}
\NormalTok{  p }\OtherTok{\textless{}{-}} \FunctionTok{nrow}\NormalTok{(W)}
\NormalTok{  h }\OtherTok{\textless{}{-}} \FunctionTok{ncol}\NormalTok{(W)}

\NormalTok{  ss }\OtherTok{\textless{}{-}} \FunctionTok{colSums}\NormalTok{(T}\SpecialCharTok{\^{}}\DecValTok{2}\NormalTok{) }\SpecialCharTok{*} \FunctionTok{colSums}\NormalTok{(Q}\SpecialCharTok{\^{}}\DecValTok{2}\NormalTok{)}
\NormalTok{  total\_ss }\OtherTok{\textless{}{-}} \FunctionTok{sum}\NormalTok{(ss)}

\NormalTok{  vip }\OtherTok{\textless{}{-}} \FunctionTok{sapply}\NormalTok{(}\FunctionTok{seq\_len}\NormalTok{(p), }\ControlFlowTok{function}\NormalTok{(j) \{}
    \FunctionTok{sqrt}\NormalTok{(p }\SpecialCharTok{*} \FunctionTok{sum}\NormalTok{(ss }\SpecialCharTok{*}\NormalTok{ (W[j, ]}\SpecialCharTok{\^{}}\DecValTok{2} \SpecialCharTok{/} \FunctionTok{colSums}\NormalTok{(W}\SpecialCharTok{\^{}}\DecValTok{2}\NormalTok{))) }\SpecialCharTok{/}\NormalTok{ total\_ss)}
\NormalTok{  \})}

  \FunctionTok{tibble}\NormalTok{(}\AttributeTok{variable =} \FunctionTok{rownames}\NormalTok{(W), }\AttributeTok{VIP =}\NormalTok{ vip) }\SpecialCharTok{\%\textgreater{}\%} \FunctionTok{arrange}\NormalTok{(}\FunctionTok{desc}\NormalTok{(VIP))}
\NormalTok{\}}
\end{Highlighting}
\end{Shaded}

\hypertarget{tecator-fat-prediction}{%
\section{6.1 Tecator fat prediction}\label{tecator-fat-prediction}}

\hypertarget{problem-statement}{%
\subsection{Problem statement}\label{problem-statement}}

The Tecator exercise uses 215 meat samples measured at 100 infrared
frequencies to predict fat content. The chapter explicitly notes that
the predictors are ordered frequencies and therefore highly correlated,
making this a good setting for PCA, PLS, and penalized linear models.
fileciteturn0file0

\hypertarget{a-load-the-data}{%
\subsection{(a) Load the data}\label{a-load-the-data}}

\begin{Shaded}
\begin{Highlighting}[]
\FunctionTok{data}\NormalTok{(tecator)}
\FunctionTok{str}\NormalTok{(absorp)}
\end{Highlighting}
\end{Shaded}

\begin{verbatim}
##  num [1:215, 1:100] 2.62 2.83 2.58 2.82 2.79 ...
\end{verbatim}

\begin{Shaded}
\begin{Highlighting}[]
\FunctionTok{str}\NormalTok{(endpoints)}
\end{Highlighting}
\end{Shaded}

\begin{verbatim}
##  num [1:215, 1:3] 60.5 46 71 72.8 58.3 44 44 69.3 61.4 61.4 ...
\end{verbatim}

\begin{Shaded}
\begin{Highlighting}[]
\NormalTok{x\_tec }\OtherTok{\textless{}{-}} \FunctionTok{as.data.frame}\NormalTok{(absorp)}
\NormalTok{y\_tec }\OtherTok{\textless{}{-}}\NormalTok{ endpoints[, }\DecValTok{2}\NormalTok{]   }\CommentTok{\# fat content}
\end{Highlighting}
\end{Shaded}

\textbf{Answer:} the predictors are the 100 absorbance measurements and
the response is percent fat.

\hypertarget{b-use-pca-to-determine-the-effective-dimension}{%
\subsection{(b) Use PCA to determine the effective
dimension}\label{b-use-pca-to-determine-the-effective-dimension}}

Because adjacent wavelengths are physically close and chemically
related, we expect strong local correlation and a much smaller effective
dimension than 100. The direct way to answer the question is to
center/scale the predictors, run PCA, and inspect the cumulative
variance curve.

\begin{Shaded}
\begin{Highlighting}[]
\NormalTok{pre\_tec }\OtherTok{\textless{}{-}} \FunctionTok{preProcess}\NormalTok{(x\_tec, }\AttributeTok{method =} \FunctionTok{c}\NormalTok{(}\StringTok{"center"}\NormalTok{, }\StringTok{"scale"}\NormalTok{))}
\NormalTok{x\_tec\_cs }\OtherTok{\textless{}{-}} \FunctionTok{predict}\NormalTok{(pre\_tec, x\_tec)}

\NormalTok{pca\_tec }\OtherTok{\textless{}{-}} \FunctionTok{prcomp}\NormalTok{(x\_tec\_cs)}
\NormalTok{var\_exp }\OtherTok{\textless{}{-}}\NormalTok{ pca\_tec}\SpecialCharTok{$}\NormalTok{sdev}\SpecialCharTok{\^{}}\DecValTok{2} \SpecialCharTok{/} \FunctionTok{sum}\NormalTok{(pca\_tec}\SpecialCharTok{$}\NormalTok{sdev}\SpecialCharTok{\^{}}\DecValTok{2}\NormalTok{)}
\NormalTok{cum\_var }\OtherTok{\textless{}{-}} \FunctionTok{cumsum}\NormalTok{(var\_exp)}

\NormalTok{eff\_dim\_90 }\OtherTok{\textless{}{-}} \FunctionTok{which}\NormalTok{(cum\_var }\SpecialCharTok{\textgreater{}=} \FloatTok{0.90}\NormalTok{)[}\DecValTok{1}\NormalTok{]}
\NormalTok{eff\_dim\_95 }\OtherTok{\textless{}{-}} \FunctionTok{which}\NormalTok{(cum\_var }\SpecialCharTok{\textgreater{}=} \FloatTok{0.95}\NormalTok{)[}\DecValTok{1}\NormalTok{]}
\NormalTok{eff\_dim\_99 }\OtherTok{\textless{}{-}} \FunctionTok{which}\NormalTok{(cum\_var }\SpecialCharTok{\textgreater{}=} \FloatTok{0.99}\NormalTok{)[}\DecValTok{1}\NormalTok{]}

\FunctionTok{data.frame}\NormalTok{(}
  \AttributeTok{threshold =} \FunctionTok{c}\NormalTok{(}\StringTok{"90\%"}\NormalTok{, }\StringTok{"95\%"}\NormalTok{, }\StringTok{"99\%"}\NormalTok{),}
  \AttributeTok{components =} \FunctionTok{c}\NormalTok{(eff\_dim\_90, eff\_dim\_95, eff\_dim\_99)}
\NormalTok{)}
\end{Highlighting}
\end{Shaded}

\begin{verbatim}
##   threshold components
## 1       90%          1
## 2       95%          1
## 3       99%          2
\end{verbatim}

\begin{Shaded}
\begin{Highlighting}[]
\FunctionTok{ggplot}\NormalTok{(}\FunctionTok{tibble}\NormalTok{(}\AttributeTok{component =} \FunctionTok{seq\_along}\NormalTok{(cum\_var), }\AttributeTok{cum\_var =}\NormalTok{ cum\_var),}
       \FunctionTok{aes}\NormalTok{(component, cum\_var)) }\SpecialCharTok{+}
  \FunctionTok{geom\_line}\NormalTok{() }\SpecialCharTok{+}
  \FunctionTok{geom\_point}\NormalTok{() }\SpecialCharTok{+}
  \FunctionTok{geom\_hline}\NormalTok{(}\AttributeTok{yintercept =} \FunctionTok{c}\NormalTok{(.}\DecValTok{90}\NormalTok{, .}\DecValTok{95}\NormalTok{), }\AttributeTok{linetype =} \DecValTok{2}\NormalTok{, }\AttributeTok{linewidth =} \FloatTok{0.3}\NormalTok{, }\AttributeTok{color =} \StringTok{"gray45"}\NormalTok{) }\SpecialCharTok{+}
  \FunctionTok{labs}\NormalTok{(}\AttributeTok{x =} \StringTok{"Principal component"}\NormalTok{, }\AttributeTok{y =} \StringTok{"Cumulative variance explained"}\NormalTok{,}
       \AttributeTok{title =} \StringTok{"Tecator PCA cumulative variance"}\NormalTok{)}
\end{Highlighting}
\end{Shaded}

\begin{center}\includegraphics{chapter6_files/figure-latex/tecator-pca-1} \end{center}

\textbf{Answer:} use the table above. For reporting, I would call the
\textbf{effective dimension} the smallest number of PCs needed to
explain about 90--95\% of the variance, since that is the practical
low-dimensional summary of the ordered spectra.

\hypertarget{c-split-the-data-preprocess-and-fit-all-chapter-6-models}{%
\subsection{(c) Split the data, preprocess, and fit all Chapter 6
models}\label{c-split-the-data-preprocess-and-fit-all-chapter-6-models}}

We will compare: - ordinary linear regression, - PLS, - ridge
regression, - lasso, - elastic net.

\begin{Shaded}
\begin{Highlighting}[]
\FunctionTok{set.seed}\NormalTok{(}\DecValTok{2026}\NormalTok{)}
\NormalTok{idx\_tec }\OtherTok{\textless{}{-}} \FunctionTok{createDataPartition}\NormalTok{(y\_tec, }\AttributeTok{p =} \FloatTok{0.80}\NormalTok{, }\AttributeTok{list =} \ConstantTok{FALSE}\NormalTok{)}

\NormalTok{x\_tec\_train }\OtherTok{\textless{}{-}}\NormalTok{ x\_tec[idx\_tec, ]}
\NormalTok{x\_tec\_test  }\OtherTok{\textless{}{-}}\NormalTok{ x\_tec[}\SpecialCharTok{{-}}\NormalTok{idx\_tec, ]}
\NormalTok{y\_tec\_train }\OtherTok{\textless{}{-}}\NormalTok{ y\_tec[idx\_tec]}
\NormalTok{y\_tec\_test  }\OtherTok{\textless{}{-}}\NormalTok{ y\_tec[}\SpecialCharTok{{-}}\NormalTok{idx\_tec]}

\NormalTok{ctrl\_reg }\OtherTok{\textless{}{-}} \FunctionTok{trainControl}\NormalTok{(}\AttributeTok{method =} \StringTok{"cv"}\NormalTok{, }\AttributeTok{number =} \DecValTok{5}\NormalTok{)}
\end{Highlighting}
\end{Shaded}

\hypertarget{linear-regression}{%
\subsubsection{Linear regression}\label{linear-regression}}

\begin{Shaded}
\begin{Highlighting}[]
\NormalTok{tec\_lm }\OtherTok{\textless{}{-}} \FunctionTok{train}\NormalTok{(}
  \AttributeTok{x =}\NormalTok{ x\_tec\_train,}
  \AttributeTok{y =}\NormalTok{ y\_tec\_train,}
  \AttributeTok{method =} \StringTok{"lm"}\NormalTok{,}
  \AttributeTok{preProcess =} \FunctionTok{c}\NormalTok{(}\StringTok{"center"}\NormalTok{, }\StringTok{"scale"}\NormalTok{),}
  \AttributeTok{trControl =}\NormalTok{ ctrl\_reg,}
  \AttributeTok{metric =} \StringTok{"RMSE"}
\NormalTok{)}
\NormalTok{tec\_lm}
\end{Highlighting}
\end{Shaded}

\begin{verbatim}
## Linear Regression 
## 
## 174 samples
## 100 predictors
## 
## Pre-processing: centered (100), scaled (100) 
## Resampling: Cross-Validated (5 fold) 
## Summary of sample sizes: 139, 141, 139, 139, 138 
## Resampling results:
## 
##   RMSE      Rsquared   MAE     
##   3.676698  0.9209796  2.474238
## 
## Tuning parameter 'intercept' was held constant at a value of TRUE
\end{verbatim}

\textbf{Answer:} linear regression has no tuning parameter, so the
``optimal'' value is not applicable.

\hypertarget{pls}{%
\subsubsection{PLS}\label{pls}}

\begin{Shaded}
\begin{Highlighting}[]
\NormalTok{pls\_grid }\OtherTok{\textless{}{-}} \FunctionTok{expand.grid}\NormalTok{(}\AttributeTok{ncomp =} \DecValTok{1}\SpecialCharTok{:}\DecValTok{15}\NormalTok{)}

\NormalTok{tec\_pls }\OtherTok{\textless{}{-}} \FunctionTok{train}\NormalTok{(}
  \AttributeTok{x =}\NormalTok{ x\_tec\_train,}
  \AttributeTok{y =}\NormalTok{ y\_tec\_train,}
  \AttributeTok{method =} \StringTok{"pls"}\NormalTok{,}
  \AttributeTok{tuneGrid =}\NormalTok{ pls\_grid,}
  \AttributeTok{preProcess =} \FunctionTok{c}\NormalTok{(}\StringTok{"center"}\NormalTok{, }\StringTok{"scale"}\NormalTok{),}
  \AttributeTok{trControl =}\NormalTok{ ctrl\_reg,}
  \AttributeTok{metric =} \StringTok{"RMSE"}
\NormalTok{)}
\NormalTok{tec\_pls}
\end{Highlighting}
\end{Shaded}

\begin{verbatim}
## Partial Least Squares 
## 
## 174 samples
## 100 predictors
## 
## Pre-processing: centered (100), scaled (100) 
## Resampling: Cross-Validated (5 fold) 
## Summary of sample sizes: 139, 139, 138, 139, 141 
## Resampling results across tuning parameters:
## 
##   ncomp  RMSE       Rsquared   MAE     
##    1     11.622198  0.2118231  9.288816
##    2      7.960089  0.6544242  6.185909
##    3      5.340941  0.8345996  4.075444
##    4      4.355952  0.8956003  3.573886
##    5      3.259243  0.9442289  2.526937
##    6      3.192745  0.9457547  2.499605
##    7      3.211157  0.9451586  2.408406
##    8      3.012660  0.9525152  2.248368
##    9      3.041014  0.9524221  2.282326
##   10      2.973371  0.9544661  2.259329
##   11      2.857732  0.9588941  2.171496
##   12      2.523195  0.9676535  1.879841
##   13      2.374429  0.9705027  1.735981
##   14      2.208817  0.9728180  1.649793
##   15      2.384897  0.9669135  1.734478
## 
## RMSE was used to select the optimal model using the smallest value.
## The final value used for the model was ncomp = 14.
\end{verbatim}

\begin{Shaded}
\begin{Highlighting}[]
\FunctionTok{plot\_tuning}\NormalTok{(tec\_pls, }\AttributeTok{x =} \StringTok{"ncomp"}\NormalTok{, }\AttributeTok{y =} \StringTok{"RMSE"}\NormalTok{, }\AttributeTok{title =} \StringTok{"Tecator PLS tuning"}\NormalTok{)}
\end{Highlighting}
\end{Shaded}

\begin{center}\includegraphics{chapter6_files/figure-latex/tecator-pls-1} \end{center}

\textbf{Answer:} the optimal number of latent components is the
\texttt{ncomp} value selected by the model summary above.

\hypertarget{ridge}{%
\subsubsection{Ridge}\label{ridge}}

\begin{Shaded}
\begin{Highlighting}[]
\NormalTok{ridge\_grid }\OtherTok{\textless{}{-}} \FunctionTok{expand.grid}\NormalTok{(}\AttributeTok{lambda =} \FunctionTok{seq}\NormalTok{(}\DecValTok{0}\NormalTok{, }\FloatTok{0.20}\NormalTok{, }\AttributeTok{length.out =} \DecValTok{15}\NormalTok{))}

\NormalTok{tec\_ridge }\OtherTok{\textless{}{-}} \FunctionTok{train}\NormalTok{(}
  \AttributeTok{x =}\NormalTok{ x\_tec\_train,}
  \AttributeTok{y =}\NormalTok{ y\_tec\_train,}
  \AttributeTok{method =} \StringTok{"ridge"}\NormalTok{,}
  \AttributeTok{tuneGrid =}\NormalTok{ ridge\_grid,}
  \AttributeTok{preProcess =} \FunctionTok{c}\NormalTok{(}\StringTok{"center"}\NormalTok{, }\StringTok{"scale"}\NormalTok{),}
  \AttributeTok{trControl =}\NormalTok{ ctrl\_reg,}
  \AttributeTok{metric =} \StringTok{"RMSE"}
\NormalTok{)}
\NormalTok{tec\_ridge}
\end{Highlighting}
\end{Shaded}

\begin{verbatim}
## Ridge Regression 
## 
## 174 samples
## 100 predictors
## 
## Pre-processing: centered (100), scaled (100) 
## Resampling: Cross-Validated (5 fold) 
## Summary of sample sizes: 139, 139, 141, 138, 139 
## Resampling results across tuning parameters:
## 
##   lambda      RMSE      Rsquared   MAE     
##   0.00000000  3.571465  0.9409278  2.377526
##   0.01428571  3.915887  0.9141461  3.249154
##   0.02857143  4.272598  0.8994477  3.483134
##   0.04285714  4.527516  0.8891768  3.628764
##   0.05714286  4.745959  0.8801071  3.751964
##   0.07142857  4.943550  0.8714022  3.880159
##   0.08571429  5.125614  0.8628019  4.014069
##   0.10000000  5.294703  0.8542244  4.135363
##   0.11428571  5.452464  0.8456501  4.243575
##   0.12857143  5.600175  0.8370830  4.340466
##   0.14285714  5.738914  0.8285348  4.429242
##   0.15714286  5.869624  0.8200199  4.510818
##   0.17142857  5.993134  0.8115523  4.585895
##   0.18571429  6.110182  0.8031452  4.657002
##   0.20000000  6.221418  0.7948102  4.726130
## 
## RMSE was used to select the optimal model using the smallest value.
## The final value used for the model was lambda = 0.
\end{verbatim}

\begin{Shaded}
\begin{Highlighting}[]
\FunctionTok{plot\_tuning}\NormalTok{(tec\_ridge, }\AttributeTok{x =} \StringTok{"lambda"}\NormalTok{, }\AttributeTok{y =} \StringTok{"RMSE"}\NormalTok{, }\AttributeTok{title =} \StringTok{"Tecator ridge tuning"}\NormalTok{)}
\end{Highlighting}
\end{Shaded}

\begin{center}\includegraphics{chapter6_files/figure-latex/tecator-ridge-1} \end{center}

\textbf{Answer:} the optimal ridge penalty is the selected
\texttt{lambda}.

\hypertarget{lasso}{%
\subsubsection{Lasso}\label{lasso}}

\begin{Shaded}
\begin{Highlighting}[]
\NormalTok{lasso\_grid }\OtherTok{\textless{}{-}} \FunctionTok{expand.grid}\NormalTok{(}
  \AttributeTok{fraction =} \FunctionTok{seq}\NormalTok{(}\FloatTok{0.10}\NormalTok{, }\FloatTok{1.00}\NormalTok{, }\AttributeTok{length.out =} \DecValTok{10}\NormalTok{),}
  \AttributeTok{lambda =} \DecValTok{0}
\NormalTok{)}

\NormalTok{tec\_lasso }\OtherTok{\textless{}{-}} \FunctionTok{train}\NormalTok{(}
  \AttributeTok{x =}\NormalTok{ x\_tec\_train,}
  \AttributeTok{y =}\NormalTok{ y\_tec\_train,}
  \AttributeTok{method =} \StringTok{"enet"}\NormalTok{,}
  \AttributeTok{tuneGrid =}\NormalTok{ lasso\_grid,}
  \AttributeTok{preProcess =} \FunctionTok{c}\NormalTok{(}\StringTok{"center"}\NormalTok{, }\StringTok{"scale"}\NormalTok{),}
  \AttributeTok{trControl =}\NormalTok{ ctrl\_reg,}
  \AttributeTok{metric =} \StringTok{"RMSE"}
\NormalTok{)}
\NormalTok{tec\_lasso}
\end{Highlighting}
\end{Shaded}

\begin{verbatim}
## Elasticnet 
## 
## 174 samples
## 100 predictors
## 
## Pre-processing: centered (100), scaled (100) 
## Resampling: Cross-Validated (5 fold) 
## Summary of sample sizes: 139, 139, 140, 140, 138 
## Resampling results across tuning parameters:
## 
##   fraction  RMSE      Rsquared   MAE     
##   0.1       2.907102  0.9431352  1.830585
##   0.2       2.884305  0.9358443  1.791299
##   0.3       2.639645  0.9511127  1.768909
##   0.4       2.556941  0.9574745  1.786130
##   0.5       2.557584  0.9596900  1.790784
##   0.6       2.662356  0.9572917  1.868958
##   0.7       2.838514  0.9530671  1.985267
##   0.8       3.016933  0.9489860  2.119754
##   0.9       3.220612  0.9438133  2.271996
##   1.0       3.455901  0.9370981  2.423820
## 
## Tuning parameter 'lambda' was held constant at a value of 0
## RMSE was used to select the optimal model using the smallest value.
## The final values used for the model were fraction = 0.4 and lambda = 0.
\end{verbatim}

\begin{Shaded}
\begin{Highlighting}[]
\FunctionTok{plot\_tuning}\NormalTok{(tec\_lasso, }\AttributeTok{x =} \StringTok{"fraction"}\NormalTok{, }\AttributeTok{y =} \StringTok{"RMSE"}\NormalTok{, }\AttributeTok{title =} \StringTok{"Tecator lasso tuning"}\NormalTok{)}
\end{Highlighting}
\end{Shaded}

\begin{center}\includegraphics{chapter6_files/figure-latex/tecator-lasso-1} \end{center}

\textbf{Answer:} because lasso is fit here as elastic net with
\texttt{lambda\ =\ 0}, the effective tuning parameter is the selected
\texttt{fraction}.

\hypertarget{elastic-net}{%
\subsubsection{Elastic net}\label{elastic-net}}

\begin{Shaded}
\begin{Highlighting}[]
\NormalTok{enet\_grid }\OtherTok{\textless{}{-}} \FunctionTok{expand.grid}\NormalTok{(}
  \AttributeTok{fraction =} \FunctionTok{seq}\NormalTok{(}\FloatTok{0.10}\NormalTok{, }\FloatTok{1.00}\NormalTok{, }\AttributeTok{length.out =} \DecValTok{10}\NormalTok{),}
  \AttributeTok{lambda =} \FunctionTok{c}\NormalTok{(}\FloatTok{0.00}\NormalTok{, }\FloatTok{0.01}\NormalTok{, }\FloatTok{0.05}\NormalTok{)}
\NormalTok{)}

\NormalTok{tec\_enet }\OtherTok{\textless{}{-}} \FunctionTok{train}\NormalTok{(}
  \AttributeTok{x =}\NormalTok{ x\_tec\_train,}
  \AttributeTok{y =}\NormalTok{ y\_tec\_train,}
  \AttributeTok{method =} \StringTok{"enet"}\NormalTok{,}
  \AttributeTok{tuneGrid =}\NormalTok{ enet\_grid,}
  \AttributeTok{preProcess =} \FunctionTok{c}\NormalTok{(}\StringTok{"center"}\NormalTok{, }\StringTok{"scale"}\NormalTok{),}
  \AttributeTok{trControl =}\NormalTok{ ctrl\_reg,}
  \AttributeTok{metric =} \StringTok{"RMSE"}
\NormalTok{)}
\NormalTok{tec\_enet}
\end{Highlighting}
\end{Shaded}

\begin{verbatim}
## Elasticnet 
## 
## 174 samples
## 100 predictors
## 
## Pre-processing: centered (100), scaled (100) 
## Resampling: Cross-Validated (5 fold) 
## Summary of sample sizes: 140, 138, 140, 139, 139 
## Resampling results across tuning parameters:
## 
##   lambda  fraction  RMSE       Rsquared   MAE     
##   0.00    0.1        2.332738  0.9703437  1.643515
##   0.00    0.2        2.311098  0.9726368  1.562452
##   0.00    0.3        2.262330  0.9739525  1.535520
##   0.00    0.4        2.214151  0.9752318  1.518391
##   0.00    0.5        2.321496  0.9730273  1.596747
##   0.00    0.6        2.509688  0.9681915  1.709508
##   0.00    0.7        2.680138  0.9629406  1.817880
##   0.00    0.8        2.970415  0.9530404  1.949865
##   0.00    0.9        3.343108  0.9377845  2.080041
##   0.00    1.0        3.803424  0.9158416  2.262305
##   0.01    0.1       10.063802  0.4602560  8.105926
##   0.01    0.2        8.529842  0.6586425  6.899677
##   0.01    0.3        7.144853  0.7749543  5.751698
##   0.01    0.4        6.072845  0.8288275  4.874107
##   0.01    0.5        5.277351  0.8543607  4.216315
##   0.01    0.6        4.657970  0.8754753  3.724478
##   0.01    0.7        4.286789  0.8889674  3.474962
##   0.01    0.8        4.091397  0.8971487  3.335886
##   0.01    0.9        3.975066  0.9027201  3.246501
##   0.01    1.0        3.908431  0.9061802  3.191396
##   0.05    0.1       10.479018  0.3933748  8.409536
##   0.05    0.2        9.327365  0.5552332  7.510074
##   0.05    0.3        8.263766  0.6766856  6.655012
##   0.05    0.4        7.316634  0.7552644  5.870763
##   0.05    0.5        6.606939  0.7956690  5.271769
##   0.05    0.6        6.061260  0.8186246  4.818116
##   0.05    0.7        5.569625  0.8380743  4.410887
##   0.05    0.8        5.175334  0.8526045  4.079815
##   0.05    0.9        4.912878  0.8612679  3.860406
##   0.05    1.0        4.757530  0.8659618  3.754580
## 
## RMSE was used to select the optimal model using the smallest value.
## The final values used for the model were fraction = 0.4 and lambda = 0.
\end{verbatim}

\begin{Shaded}
\begin{Highlighting}[]
\FunctionTok{plot\_tuning}\NormalTok{(tec\_enet, }\AttributeTok{x =} \StringTok{"fraction"}\NormalTok{, }\AttributeTok{y =} \StringTok{"RMSE"}\NormalTok{, }\AttributeTok{color =} \StringTok{"lambda"}\NormalTok{,}
            \AttributeTok{title =} \StringTok{"Tecator elastic net tuning"}\NormalTok{)}
\end{Highlighting}
\end{Shaded}

\begin{center}\includegraphics{chapter6_files/figure-latex/tecator-enet-1} \end{center}

\textbf{Answer:} the optimal elastic-net tuning parameters are the
selected \texttt{lambda} and \texttt{fraction}.

\hypertarget{d-which-model-has-the-best-predictive-ability}{%
\subsection{(d) Which model has the best predictive
ability?}\label{d-which-model-has-the-best-predictive-ability}}

\begin{Shaded}
\begin{Highlighting}[]
\NormalTok{tec\_models }\OtherTok{\textless{}{-}} \FunctionTok{list}\NormalTok{(}
  \AttributeTok{LM =}\NormalTok{ tec\_lm,}
  \AttributeTok{PLS =}\NormalTok{ tec\_pls,}
  \AttributeTok{Ridge =}\NormalTok{ tec\_ridge,}
  \AttributeTok{Lasso =}\NormalTok{ tec\_lasso,}
  \AttributeTok{ElasticNet =}\NormalTok{ tec\_enet}
\NormalTok{)}

\NormalTok{tec\_test\_perf }\OtherTok{\textless{}{-}} \FunctionTok{bind\_rows}\NormalTok{(}\FunctionTok{lapply}\NormalTok{(}\FunctionTok{names}\NormalTok{(tec\_models), }\ControlFlowTok{function}\NormalTok{(nm) \{}
\NormalTok{  pred }\OtherTok{\textless{}{-}} \FunctionTok{predict}\NormalTok{(tec\_models[[nm]], x\_tec\_test)}
  \FunctionTok{tibble}\NormalTok{(}\AttributeTok{model =}\NormalTok{ nm,}
         \AttributeTok{RMSE =}\NormalTok{ caret}\SpecialCharTok{::}\FunctionTok{RMSE}\NormalTok{(pred, y\_tec\_test),}
         \AttributeTok{Rsquared =}\NormalTok{ caret}\SpecialCharTok{::}\FunctionTok{R2}\NormalTok{(pred, y\_tec\_test))}
\NormalTok{\})) }\SpecialCharTok{\%\textgreater{}\%} \FunctionTok{arrange}\NormalTok{(RMSE)}

\NormalTok{tec\_test\_perf}
\end{Highlighting}
\end{Shaded}

\begin{verbatim}
## # A tibble: 5 x 3
##   model       RMSE Rsquared
##   <chr>      <dbl>    <dbl>
## 1 PLS         2.40    0.959
## 2 Lasso       4.69    0.861
## 3 ElasticNet  4.69    0.861
## 4 Ridge       5.68    0.803
## 5 LM          5.68    0.803
\end{verbatim}

To compare resampled distributions:

\begin{Shaded}
\begin{Highlighting}[]
\NormalTok{tec\_resamp }\OtherTok{\textless{}{-}} \FunctionTok{resamples}\NormalTok{(tec\_models)}
\FunctionTok{summary}\NormalTok{(tec\_resamp)}
\end{Highlighting}
\end{Shaded}

\begin{verbatim}
## 
## Call:
## summary.resamples(object = tec_resamp)
## 
## Models: LM, PLS, Ridge, Lasso, ElasticNet 
## Number of resamples: 5 
## 
## MAE 
##                Min.  1st Qu.   Median     Mean  3rd Qu.     Max. NA's
## LM         1.682484 2.116087 2.161943 2.474238 3.083070 3.327607    0
## PLS        1.533240 1.572268 1.644400 1.649793 1.649643 1.849416    0
## Ridge      1.763377 2.407593 2.407842 2.377526 2.487197 2.821621    0
## Lasso      1.278178 1.398628 1.482443 1.786130 1.717621 3.053780    0
## ElasticNet 1.329788 1.339883 1.371867 1.518391 1.662829 1.887586    0
## 
## RMSE 
##                Min.  1st Qu.   Median     Mean  3rd Qu.     Max. NA's
## LM         2.432179 2.697501 3.036375 3.676698 4.922797 5.294640    0
## PLS        1.953498 2.023573 2.195984 2.208817 2.289804 2.581229    0
## Ridge      2.609326 3.448244 3.807463 3.571465 3.979388 4.012902    0
## Lasso      1.770939 1.949239 1.997812 2.556941 2.357243 4.709474    0
## ElasticNet 1.803522 1.897477 1.959831 2.214151 2.627699 2.782224    0
## 
## Rsquared 
##                 Min.   1st Qu.    Median      Mean   3rd Qu.      Max. NA's
## LM         0.8607977 0.8699964 0.9418633 0.9209796 0.9645767 0.9676640    0
## PLS        0.9693971 0.9704690 0.9708953 0.9728180 0.9755778 0.9777508    0
## Ridge      0.9233091 0.9283660 0.9421674 0.9409278 0.9515672 0.9592293    0
## Lasso      0.8749808 0.9717348 0.9786445 0.9574745 0.9795798 0.9824324    0
## ElasticNet 0.9646928 0.9749136 0.9750975 0.9752318 0.9761251 0.9853301    0
\end{verbatim}

\begin{Shaded}
\begin{Highlighting}[]
\FunctionTok{summary}\NormalTok{(}\FunctionTok{diff}\NormalTok{(tec\_resamp))}
\end{Highlighting}
\end{Shaded}

\begin{verbatim}
## 
## Call:
## summary.diff.resamples(object = diff(tec_resamp))
## 
## p-value adjustment: bonferroni 
## Upper diagonal: estimates of the difference
## Lower diagonal: p-value for H0: difference = 0
## 
## MAE 
##            LM     PLS      Ridge    Lasso    ElasticNet
## LM                 0.82444  0.09671  0.68811  0.95585  
## PLS        0.5242          -0.72773 -0.13634  0.13140  
## Ridge      1.0000 0.1790             0.59140  0.85914  
## Lasso      0.5457 1.0000   0.8436             0.26774  
## ElasticNet 0.3335 1.0000   0.1225   1.0000             
## 
## RMSE 
##            LM      PLS       Ridge     Lasso     ElasticNet
## LM                  1.467881  0.105234  1.119757  1.462548 
## PLS        0.63197           -1.362647 -0.348124 -0.005333 
## Ridge      1.00000 0.16583              1.014523  1.357314 
## Lasso      0.99476 1.00000   1.00000              0.342791 
## ElasticNet 0.71075 1.00000   0.09737   1.00000             
## 
## Rsquared 
##            LM      PLS       Ridge     Lasso     ElasticNet
## LM                 -0.051838 -0.019948 -0.036495 -0.054252 
## PLS        0.76936            0.031890  0.015344 -0.002414 
## Ridge      1.00000 0.10855             -0.016547 -0.034304 
## Lasso      1.00000 1.00000   1.00000             -0.017757 
## ElasticNet 0.66070 1.00000   0.05633   1.00000
\end{verbatim}

\textbf{Answer:} the best predictive model is the one with the
\textbf{lowest test-set RMSE} and, secondarily, the highest test-set
\(R^2\). If the paired resampling comparison shows very small
differences, then I would prefer the simpler or more stable model among
the top performers.

\hypertarget{e-which-model-would-i-use}{%
\subsection{(e) Which model would I
use?}\label{e-which-model-would-i-use}}

\textbf{Direct answer:} for Tecator I would usually choose \textbf{PLS}
unless ridge or elastic net is clearly better on the test set.

\textbf{Why:} 1. The spectra are extremely correlated and ordered, which
is exactly where latent-component models shine. 2. PLS uses the response
while constructing components, which is preferable to unsupervised
PCA/PCR in this setting. Chapter 6 explicitly recommends PLS over PCR
when correlated predictors are present and a linear-regression-type
solution is desired. fileciteturn3file8L14-L21 3. PLS gives good
predictive performance with fewer effective dimensions and a coherent
scientific story: the prediction is driven by a small number of latent
spectral directions.

If ridge or elastic net gives materially smaller test RMSE, I would use
that instead for pure prediction. If performance is essentially tied, I
would still pick PLS.

\hypertarget{permeability}{%
\section{6.2 Permeability}\label{permeability}}

\hypertarget{problem-statement-1}{%
\subsection{Problem statement}\label{problem-statement-1}}

The permeability exercise uses 1,107 binary molecular fingerprint
predictors on only 165 compounds. The book describes this data set as a
classic high-dimensional problem with sparse predictors and a skewed
response, making it a strong candidate for filtering and latent-variable
modeling. fileciteturn4file0L24-L26

\hypertarget{a-load-the-data-1}{%
\subsection{(a) Load the data}\label{a-load-the-data-1}}

\begin{Shaded}
\begin{Highlighting}[]
\FunctionTok{data}\NormalTok{(permeability)}
\FunctionTok{str}\NormalTok{(fingerprints)}
\end{Highlighting}
\end{Shaded}

\begin{verbatim}
##  num [1:165, 1:1107] 0 0 0 0 0 0 0 0 0 0 ...
##  - attr(*, "dimnames")=List of 2
##   ..$ : chr [1:165] "1" "2" "3" "4" ...
##   ..$ : chr [1:1107] "X1" "X2" "X3" "X4" ...
\end{verbatim}

\begin{Shaded}
\begin{Highlighting}[]
\FunctionTok{str}\NormalTok{(permeability)}
\end{Highlighting}
\end{Shaded}

\begin{verbatim}
##  num [1:165, 1] 12.52 1.12 19.41 1.73 1.68 ...
##  - attr(*, "dimnames")=List of 2
##   ..$ : chr [1:165] "1" "2" "3" "4" ...
##   ..$ : chr "permeability"
\end{verbatim}

\begin{Shaded}
\begin{Highlighting}[]
\NormalTok{x\_perm }\OtherTok{\textless{}{-}} \FunctionTok{as.data.frame}\NormalTok{(fingerprints)}
\NormalTok{y\_perm }\OtherTok{\textless{}{-}}\NormalTok{ permeability}
\end{Highlighting}
\end{Shaded}

\hypertarget{b-remove-very-sparse-predictors-with-nearzerovar}{%
\subsection{\texorpdfstring{(b) Remove very sparse predictors with
\texttt{nearZeroVar}}{(b) Remove very sparse predictors with nearZeroVar}}\label{b-remove-very-sparse-predictors-with-nearzerovar}}

\begin{Shaded}
\begin{Highlighting}[]
\NormalTok{nzv\_idx }\OtherTok{\textless{}{-}} \FunctionTok{nearZeroVar}\NormalTok{(x\_perm)}
\FunctionTok{length}\NormalTok{(nzv\_idx)}
\end{Highlighting}
\end{Shaded}

\begin{verbatim}
## [1] 719
\end{verbatim}

\begin{Shaded}
\begin{Highlighting}[]
\NormalTok{x\_perm\_filt }\OtherTok{\textless{}{-}}\NormalTok{ x\_perm[, }\SpecialCharTok{{-}}\NormalTok{nzv\_idx]}
\FunctionTok{dim}\NormalTok{(x\_perm\_filt)}
\end{Highlighting}
\end{Shaded}

\begin{verbatim}
## [1] 165 388
\end{verbatim}

\textbf{Answer:} the number of predictors left for modeling is
\texttt{ncol(x\_perm\_filt)}.

\hypertarget{c-split-preprocess-and-tune-a-pls-model}{%
\subsection{(c) Split, preprocess, and tune a PLS
model}\label{c-split-preprocess-and-tune-a-pls-model}}

Because the data are high-dimensional (\(p \gg n\)), ordinary least
squares is not an attractive primary model. Chapter 6 emphasizes that
PLS is a strong alternative in exactly this kind of situation.
fileciteturn3file8L1-L21

\begin{Shaded}
\begin{Highlighting}[]
\FunctionTok{set.seed}\NormalTok{(}\DecValTok{2026}\NormalTok{)}
\NormalTok{idx\_perm }\OtherTok{\textless{}{-}} \FunctionTok{createDataPartition}\NormalTok{(y\_perm, }\AttributeTok{p =} \FloatTok{0.80}\NormalTok{, }\AttributeTok{list =} \ConstantTok{FALSE}\NormalTok{)}

\NormalTok{x\_perm\_train }\OtherTok{\textless{}{-}}\NormalTok{ x\_perm\_filt[idx\_perm, ]}
\NormalTok{x\_perm\_test  }\OtherTok{\textless{}{-}}\NormalTok{ x\_perm\_filt[}\SpecialCharTok{{-}}\NormalTok{idx\_perm, ]}
\NormalTok{y\_perm\_train }\OtherTok{\textless{}{-}}\NormalTok{ y\_perm[idx\_perm]}
\NormalTok{y\_perm\_test  }\OtherTok{\textless{}{-}}\NormalTok{ y\_perm[}\SpecialCharTok{{-}}\NormalTok{idx\_perm]}

\NormalTok{perm\_ctrl }\OtherTok{\textless{}{-}} \FunctionTok{trainControl}\NormalTok{(}\AttributeTok{method =} \StringTok{"cv"}\NormalTok{, }\AttributeTok{number =} \DecValTok{5}\NormalTok{)}
\NormalTok{perm\_grid }\OtherTok{\textless{}{-}} \FunctionTok{expand.grid}\NormalTok{(}\AttributeTok{ncomp =} \DecValTok{1}\SpecialCharTok{:}\FunctionTok{min}\NormalTok{(}\DecValTok{15}\NormalTok{, }\FunctionTok{ncol}\NormalTok{(x\_perm\_train)))}

\NormalTok{perm\_pls }\OtherTok{\textless{}{-}} \FunctionTok{train}\NormalTok{(}
  \AttributeTok{x =}\NormalTok{ x\_perm\_train,}
  \AttributeTok{y =}\NormalTok{ y\_perm\_train,}
  \AttributeTok{method =} \StringTok{"pls"}\NormalTok{,}
  \AttributeTok{tuneGrid =}\NormalTok{ perm\_grid,}
  \AttributeTok{preProcess =} \FunctionTok{c}\NormalTok{(}\StringTok{"center"}\NormalTok{, }\StringTok{"scale"}\NormalTok{),}
  \AttributeTok{trControl =}\NormalTok{ perm\_ctrl,}
  \AttributeTok{metric =} \StringTok{"Rsquared"}
\NormalTok{)}
\NormalTok{perm\_pls}
\end{Highlighting}
\end{Shaded}

\begin{verbatim}
## Partial Least Squares 
## 
## 133 samples
## 388 predictors
## 
## Pre-processing: centered (388), scaled (388) 
## Resampling: Cross-Validated (5 fold) 
## Summary of sample sizes: 106, 105, 108, 106, 107 
## Resampling results across tuning parameters:
## 
##   ncomp  RMSE      Rsquared   MAE      
##    1     13.43505  0.2928362  10.191173
##    2     12.20378  0.4383808   8.704919
##    3     12.22274  0.4420488   9.176851
##    4     12.33205  0.4533712   9.439310
##    5     11.99808  0.4911597   8.915141
##    6     11.71345  0.5021128   8.774153
##    7     11.63672  0.5075570   8.811792
##    8     12.10665  0.4755174   9.226860
##    9     12.36300  0.4655705   9.255613
##   10     12.58343  0.4642721   9.565218
##   11     13.11168  0.4364529   9.781295
##   12     13.37359  0.4232621   9.956409
##   13     13.68520  0.4092271  10.001258
##   14     14.14268  0.3839722  10.299231
##   15     14.50802  0.3603010  10.618110
## 
## Rsquared was used to select the optimal model using the largest value.
## The final value used for the model was ncomp = 7.
\end{verbatim}

\begin{Shaded}
\begin{Highlighting}[]
\FunctionTok{plot\_tuning}\NormalTok{(perm\_pls, }\AttributeTok{x =} \StringTok{"ncomp"}\NormalTok{, }\AttributeTok{y =} \StringTok{"Rsquared"}\NormalTok{, }\AttributeTok{title =} \StringTok{"Permeability PLS tuning"}\NormalTok{)}
\end{Highlighting}
\end{Shaded}

\begin{center}\includegraphics{chapter6_files/figure-latex/perm-split-pls-1} \end{center}

\textbf{Answer:} the optimal number of latent variables is the selected
\texttt{ncomp}, and the resampled estimate of \(R^2\) is the
corresponding \texttt{Rsquared} value in the output.

\hypertarget{d-predict-the-test-set-and-estimate-r2}{%
\subsection{\texorpdfstring{(d) Predict the test set and estimate
\(R^2\)}{(d) Predict the test set and estimate R\^{}2}}\label{d-predict-the-test-set-and-estimate-r2}}

\begin{Shaded}
\begin{Highlighting}[]
\NormalTok{perm\_pls\_pred }\OtherTok{\textless{}{-}} \FunctionTok{predict}\NormalTok{(perm\_pls, x\_perm\_test)}
\FunctionTok{rmse\_r2}\NormalTok{(y\_perm\_test, perm\_pls\_pred)}
\end{Highlighting}
\end{Shaded}

\begin{verbatim}
##       RMSE   Rsquared        MAE 
## 11.1887496  0.4814568  8.4030025
\end{verbatim}

\begin{Shaded}
\begin{Highlighting}[]
\FunctionTok{plot\_obs\_pred}\NormalTok{(y\_perm\_test, perm\_pls\_pred, }\StringTok{"Permeability: PLS on test set"}\NormalTok{)}
\end{Highlighting}
\end{Shaded}

\begin{center}\includegraphics{chapter6_files/figure-latex/perm-test-pls-1} \end{center}

\textbf{Answer:} the test-set estimate of \(R^2\) is the
\texttt{Rsquared} value printed above.

\hypertarget{e-try-other-chapter-6-models}{%
\subsection{(e) Try other Chapter 6
models}\label{e-try-other-chapter-6-models}}

We will compare PLS with ridge, lasso, and elastic net. A plain linear
model is not suitable here because the original problem has far more
predictors than samples and even after filtering, instability is still
likely.

\begin{Shaded}
\begin{Highlighting}[]
\NormalTok{perm\_ridge }\OtherTok{\textless{}{-}} \FunctionTok{train}\NormalTok{(}
  \AttributeTok{x =}\NormalTok{ x\_perm\_train,}
  \AttributeTok{y =}\NormalTok{ y\_perm\_train,}
  \AttributeTok{method =} \StringTok{"ridge"}\NormalTok{,}
  \AttributeTok{tuneGrid =} \FunctionTok{expand.grid}\NormalTok{(}\AttributeTok{lambda =} \FunctionTok{seq}\NormalTok{(}\DecValTok{0}\NormalTok{, }\FloatTok{0.20}\NormalTok{, }\AttributeTok{length.out =} \DecValTok{15}\NormalTok{)),}
  \AttributeTok{preProcess =} \FunctionTok{c}\NormalTok{(}\StringTok{"center"}\NormalTok{, }\StringTok{"scale"}\NormalTok{),}
  \AttributeTok{trControl =}\NormalTok{ perm\_ctrl,}
  \AttributeTok{metric =} \StringTok{"Rsquared"}
\NormalTok{)}

\NormalTok{perm\_lasso }\OtherTok{\textless{}{-}} \FunctionTok{train}\NormalTok{(}
  \AttributeTok{x =}\NormalTok{ x\_perm\_train,}
  \AttributeTok{y =}\NormalTok{ y\_perm\_train,}
  \AttributeTok{method =} \StringTok{"enet"}\NormalTok{,}
  \AttributeTok{tuneGrid =} \FunctionTok{expand.grid}\NormalTok{(}\AttributeTok{lambda =} \DecValTok{0}\NormalTok{, }\AttributeTok{fraction =} \FunctionTok{seq}\NormalTok{(}\FloatTok{0.10}\NormalTok{, }\DecValTok{1}\NormalTok{, }\AttributeTok{length.out =} \DecValTok{10}\NormalTok{)),}
  \AttributeTok{preProcess =} \FunctionTok{c}\NormalTok{(}\StringTok{"center"}\NormalTok{, }\StringTok{"scale"}\NormalTok{),}
  \AttributeTok{trControl =}\NormalTok{ perm\_ctrl,}
  \AttributeTok{metric =} \StringTok{"Rsquared"}
\NormalTok{)}

\NormalTok{perm\_enet }\OtherTok{\textless{}{-}} \FunctionTok{train}\NormalTok{(}
  \AttributeTok{x =}\NormalTok{ x\_perm\_train,}
  \AttributeTok{y =}\NormalTok{ y\_perm\_train,}
  \AttributeTok{method =} \StringTok{"enet"}\NormalTok{,}
  \AttributeTok{tuneGrid =} \FunctionTok{expand.grid}\NormalTok{(}\AttributeTok{lambda =} \FunctionTok{c}\NormalTok{(}\FloatTok{0.00}\NormalTok{, }\FloatTok{0.01}\NormalTok{, }\FloatTok{0.05}\NormalTok{), }\AttributeTok{fraction =} \FunctionTok{seq}\NormalTok{(}\FloatTok{0.10}\NormalTok{, }\DecValTok{1}\NormalTok{, }\AttributeTok{length.out =} \DecValTok{10}\NormalTok{)),}
  \AttributeTok{preProcess =} \FunctionTok{c}\NormalTok{(}\StringTok{"center"}\NormalTok{, }\StringTok{"scale"}\NormalTok{),}
  \AttributeTok{trControl =}\NormalTok{ perm\_ctrl,}
  \AttributeTok{metric =} \StringTok{"Rsquared"}
\NormalTok{)}

\NormalTok{perm\_models }\OtherTok{\textless{}{-}} \FunctionTok{list}\NormalTok{(}\AttributeTok{PLS =}\NormalTok{ perm\_pls, }\AttributeTok{Ridge =}\NormalTok{ perm\_ridge, }\AttributeTok{Lasso =}\NormalTok{ perm\_lasso, }\AttributeTok{ElasticNet =}\NormalTok{ perm\_enet)}

\NormalTok{perm\_test\_perf }\OtherTok{\textless{}{-}} \FunctionTok{bind\_rows}\NormalTok{(}\FunctionTok{lapply}\NormalTok{(}\FunctionTok{names}\NormalTok{(perm\_models), }\ControlFlowTok{function}\NormalTok{(nm) \{}
\NormalTok{  pred }\OtherTok{\textless{}{-}} \FunctionTok{predict}\NormalTok{(perm\_models[[nm]], x\_perm\_test)}
  \FunctionTok{tibble}\NormalTok{(}\AttributeTok{model =}\NormalTok{ nm,}
         \AttributeTok{RMSE =}\NormalTok{ caret}\SpecialCharTok{::}\FunctionTok{RMSE}\NormalTok{(pred, y\_perm\_test),}
         \AttributeTok{Rsquared =}\NormalTok{ caret}\SpecialCharTok{::}\FunctionTok{R2}\NormalTok{(pred, y\_perm\_test))}
\NormalTok{\})) }\SpecialCharTok{\%\textgreater{}\%} \FunctionTok{arrange}\NormalTok{(RMSE)}

\NormalTok{perm\_test\_perf}
\end{Highlighting}
\end{Shaded}

\begin{verbatim}
## # A tibble: 4 x 3
##   model       RMSE Rsquared
##   <chr>      <dbl>    <dbl>
## 1 Lasso       9.96    0.532
## 2 Ridge      10.3     0.560
## 3 ElasticNet 10.8     0.458
## 4 PLS        11.2     0.481
\end{verbatim}

\begin{Shaded}
\begin{Highlighting}[]
\FunctionTok{summary}\NormalTok{(}\FunctionTok{resamples}\NormalTok{(perm\_models))}
\end{Highlighting}
\end{Shaded}

\begin{verbatim}
## 
## Call:
## summary.resamples(object = resamples(perm_models))
## 
## Models: PLS, Ridge, Lasso, ElasticNet 
## Number of resamples: 5 
## 
## MAE 
##                Min.  1st Qu.    Median      Mean   3rd Qu.      Max. NA's
## PLS        6.219190 7.605707  8.978883  8.811792 10.220910 11.034269    0
## Ridge      8.950801 9.743070 10.128185 10.006878 10.348320 10.864012    0
## Lasso      7.621820 8.004644  8.331757  8.598963  8.926075 10.110517    1
## ElasticNet 7.761840 8.038931  8.456838  8.527141  9.182913  9.195185    0
## 
## RMSE 
##                 Min.  1st Qu.   Median     Mean  3rd Qu.     Max. NA's
## PLS         9.039098 10.97091 11.97316 11.63672 12.75372 13.44673    0
## Ridge      11.676698 13.15692 13.66472 13.57347 14.51572 14.85329    0
## Lasso      10.026190 10.15836 10.52784 11.10009 11.46958 13.31849    1
## ElasticNet 10.622543 10.65748 11.10502 11.37240 12.08972 12.38721    0
## 
## Rsquared 
##                 Min.   1st Qu.    Median      Mean   3rd Qu.      Max. NA's
## PLS        0.4208595 0.4263345 0.4754682 0.5075570 0.5991815 0.6159415    0
## Ridge      0.3451141 0.3951728 0.4934692 0.4731071 0.5186481 0.6131314    0
## Lasso      0.4645855 0.4687302 0.4840597 0.5285329 0.5438623 0.6814267    1
## ElasticNet 0.4136885 0.4711638 0.4842472 0.5362397 0.5986140 0.7134848    0
\end{verbatim}

\begin{Shaded}
\begin{Highlighting}[]
\FunctionTok{summary}\NormalTok{(}\FunctionTok{diff}\NormalTok{(}\FunctionTok{resamples}\NormalTok{(perm\_models)))}
\end{Highlighting}
\end{Shaded}

\begin{verbatim}
## 
## Call:
## summary.diff.resamples(object = diff(resamples(perm_models)))
## 
## p-value adjustment: bonferroni 
## Upper diagonal: estimates of the difference
## Lower diagonal: p-value for H0: difference = 0
## 
## MAE 
##            PLS    Ridge   Lasso   ElasticNet
## PLS               -1.1951  0.8610  0.2847   
## Ridge      0.8884          1.6719  1.4797   
## Lasso      1.0000 0.5741          -0.1195   
## ElasticNet 1.0000 0.0210  1.0000            
## 
## RMSE 
##            PLS      Ridge    Lasso    ElasticNet
## PLS                 -1.9367   1.1860   0.2643   
## Ridge      0.181762           2.9476   2.2011   
## Lasso      1.000000 0.281484          -0.4598   
## ElasticNet 1.000000 0.009394 1.000000           
## 
## Rsquared 
##            PLS Ridge    Lasso    ElasticNet
## PLS             0.03445 -0.04388 -0.02868  
## Ridge      1            -0.09043 -0.06313  
## Lasso      1   1                 -0.02398  
## ElasticNet 1   1        1
\end{verbatim}

\textbf{Answer:} whichever model is at the top of
\texttt{perm\_test\_perf} has the best observed predictive performance.
In many high-dimensional chemistry problems, \textbf{PLS} or
\textbf{elastic net} tends to do very well.

\hypertarget{f-would-i-recommend-replacing-the-experiment}{%
\subsection{(f) Would I recommend replacing the
experiment?}\label{f-would-i-recommend-replacing-the-experiment}}

\textbf{Direct answer:} probably \textbf{not as a full replacement}
unless the test-set performance is extremely strong and the prediction
error is small enough for the business/scientific decision.

I would usually recommend the best model as a \textbf{screening tool},
not a total replacement, for three reasons: 1. The laboratory assay is
the ground truth. 2. The sample size is modest. 3. A false positive or
false negative in permeability can be expensive downstream.

So my recommendation would be: - use the model to \textbf{prioritize}
compounds, - send the most promising molecules to the assay, - keep
collecting new assay data and periodically retrain the model.

\hypertarget{chemical-manufacturing-process}{%
\section{6.3 Chemical manufacturing
process}\label{chemical-manufacturing-process}}

\hypertarget{problem-statement-2}{%
\subsection{Problem statement}\label{problem-statement-2}}

This exercise studies 57 predictors (12 biological/raw-material
variables and 45 process variables) for 176 manufacturing runs, with
percent yield as the outcome. The key applied question is not only
prediction but also which variables are actionable for improving yield.
The exercise text makes that distinction explicit.
fileciteturn2file14L1-L20

\hypertarget{a-load-the-data-2}{%
\subsection{(a) Load the data}\label{a-load-the-data-2}}

\begin{Shaded}
\begin{Highlighting}[]
\FunctionTok{data}\NormalTok{(}\StringTok{"ChemicalManufacturingProcess"}\NormalTok{)}
\FunctionTok{str}\NormalTok{(ChemicalManufacturingProcess)}
\end{Highlighting}
\end{Shaded}

\begin{verbatim}
## 'data.frame':    176 obs. of  58 variables:
##  $ Yield                 : num  38 42.4 42 41.4 42.5 ...
##  $ BiologicalMaterial01  : num  6.25 8.01 8.01 8.01 7.47 6.12 7.48 6.94 6.94 6.94 ...
##  $ BiologicalMaterial02  : num  49.6 61 61 61 63.3 ...
##  $ BiologicalMaterial03  : num  57 67.5 67.5 67.5 72.2 ...
##  $ BiologicalMaterial04  : num  12.7 14.7 14.7 14.7 14 ...
##  $ BiologicalMaterial05  : num  19.5 19.4 19.4 19.4 17.9 ...
##  $ BiologicalMaterial06  : num  43.7 53.1 53.1 53.1 54.7 ...
##  $ BiologicalMaterial07  : num  100 100 100 100 100 100 100 100 100 100 ...
##  $ BiologicalMaterial08  : num  16.7 19 19 19 18.2 ...
##  $ BiologicalMaterial09  : num  11.4 12.6 12.6 12.6 12.8 ...
##  $ BiologicalMaterial10  : num  3.46 3.46 3.46 3.46 3.05 3.78 3.04 3.85 3.85 3.85 ...
##  $ BiologicalMaterial11  : num  138 154 154 154 148 ...
##  $ BiologicalMaterial12  : num  18.8 21.1 21.1 21.1 21.1 ...
##  $ ManufacturingProcess01: num  NA 0 0 0 10.7 12 11.5 12 12 12 ...
##  $ ManufacturingProcess02: num  NA 0 0 0 0 0 0 0 0 0 ...
##  $ ManufacturingProcess03: num  NA NA NA NA NA NA 1.56 1.55 1.56 1.55 ...
##  $ ManufacturingProcess04: num  NA 917 912 911 918 924 933 929 928 938 ...
##  $ ManufacturingProcess05: num  NA 1032 1004 1015 1028 ...
##  $ ManufacturingProcess06: num  NA 210 207 213 206 ...
##  $ ManufacturingProcess07: num  NA 177 178 177 178 178 177 178 177 177 ...
##  $ ManufacturingProcess08: num  NA 178 178 177 178 178 178 178 177 177 ...
##  $ ManufacturingProcess09: num  43 46.6 45.1 44.9 45 ...
##  $ ManufacturingProcess10: num  NA NA NA NA NA NA 11.6 10.2 9.7 10.1 ...
##  $ ManufacturingProcess11: num  NA NA NA NA NA NA 11.5 11.3 11.1 10.2 ...
##  $ ManufacturingProcess12: num  NA 0 0 0 0 0 0 0 0 0 ...
##  $ ManufacturingProcess13: num  35.5 34 34.8 34.8 34.6 34 32.4 33.6 33.9 34.3 ...
##  $ ManufacturingProcess14: num  4898 4869 4878 4897 4992 ...
##  $ ManufacturingProcess15: num  6108 6095 6087 6102 6233 ...
##  $ ManufacturingProcess16: num  4682 4617 4617 4635 4733 ...
##  $ ManufacturingProcess17: num  35.5 34 34.8 34.8 33.9 33.4 33.8 33.6 33.9 35.3 ...
##  $ ManufacturingProcess18: num  4865 4867 4877 4872 4886 ...
##  $ ManufacturingProcess19: num  6049 6097 6078 6073 6102 ...
##  $ ManufacturingProcess20: num  4665 4621 4621 4611 4659 ...
##  $ ManufacturingProcess21: num  0 0 0 0 -0.7 -0.6 1.4 0 0 1 ...
##  $ ManufacturingProcess22: num  NA 3 4 5 8 9 1 2 3 4 ...
##  $ ManufacturingProcess23: num  NA 0 1 2 4 1 1 2 3 1 ...
##  $ ManufacturingProcess24: num  NA 3 4 5 18 1 1 2 3 4 ...
##  $ ManufacturingProcess25: num  4873 4869 4897 4892 4930 ...
##  $ ManufacturingProcess26: num  6074 6107 6116 6111 6151 ...
##  $ ManufacturingProcess27: num  4685 4630 4637 4630 4684 ...
##  $ ManufacturingProcess28: num  10.7 11.2 11.1 11.1 11.3 11.4 11.2 11.1 11.3 11.4 ...
##  $ ManufacturingProcess29: num  21 21.4 21.3 21.3 21.6 21.7 21.2 21.2 21.5 21.7 ...
##  $ ManufacturingProcess30: num  9.9 9.9 9.4 9.4 9 10.1 11.2 10.9 10.5 9.8 ...
##  $ ManufacturingProcess31: num  69.1 68.7 69.3 69.3 69.4 68.2 67.6 67.9 68 68.5 ...
##  $ ManufacturingProcess32: num  156 169 173 171 171 173 159 161 160 164 ...
##  $ ManufacturingProcess33: num  66 66 66 68 70 70 65 65 65 66 ...
##  $ ManufacturingProcess34: num  2.4 2.6 2.6 2.5 2.5 2.5 2.5 2.5 2.5 2.5 ...
##  $ ManufacturingProcess35: num  486 508 509 496 468 490 475 478 491 488 ...
##  $ ManufacturingProcess36: num  0.019 0.019 0.018 0.018 0.017 0.018 0.019 0.019 0.019 0.019 ...
##  $ ManufacturingProcess37: num  0.5 2 0.7 1.2 0.2 0.4 0.8 1 1.2 1.8 ...
##  $ ManufacturingProcess38: num  3 2 2 2 2 2 2 2 3 3 ...
##  $ ManufacturingProcess39: num  7.2 7.2 7.2 7.2 7.3 7.2 7.3 7.3 7.4 7.1 ...
##  $ ManufacturingProcess40: num  NA 0.1 0 0 0 0 0 0 0 0 ...
##  $ ManufacturingProcess41: num  NA 0.15 0 0 0 0 0 0 0 0 ...
##  $ ManufacturingProcess42: num  11.6 11.1 12 10.6 11 11.5 11.7 11.4 11.4 11.3 ...
##  $ ManufacturingProcess43: num  3 0.9 1 1.1 1.1 2.2 0.7 0.8 0.9 0.8 ...
##  $ ManufacturingProcess44: num  1.8 1.9 1.8 1.8 1.7 1.8 2 2 1.9 1.9 ...
##  $ ManufacturingProcess45: num  2.4 2.2 2.3 2.1 2.1 2 2.2 2.2 2.1 2.4 ...
\end{verbatim}

\begin{Shaded}
\begin{Highlighting}[]
\NormalTok{x\_chem }\OtherTok{\textless{}{-}}\NormalTok{ ChemicalManufacturingProcess }\SpecialCharTok{\%\textgreater{}\%} \FunctionTok{select}\NormalTok{(}\SpecialCharTok{{-}}\NormalTok{Yield)}
\NormalTok{y\_chem }\OtherTok{\textless{}{-}}\NormalTok{ ChemicalManufacturingProcess}\SpecialCharTok{$}\NormalTok{Yield}
\end{Highlighting}
\end{Shaded}

\hypertarget{b-impute-missing-values}{%
\subsection{(b) Impute missing values}\label{b-impute-missing-values}}

The problem explicitly tells us that a small percentage of predictor
cells are missing. The simplest reasonable Chapter 6-compatible workflow
is median imputation followed by centering/scaling, though
k-nearest-neighbor imputation is also fine. Because the data set is
small, I prefer \textbf{kNN imputation}.

\begin{Shaded}
\begin{Highlighting}[]
\FunctionTok{colSums}\NormalTok{(}\FunctionTok{is.na}\NormalTok{(x\_chem)) }\SpecialCharTok{\%\textgreater{}\%} \FunctionTok{sort}\NormalTok{(}\AttributeTok{decreasing =} \ConstantTok{TRUE}\NormalTok{) }\SpecialCharTok{\%\textgreater{}\%} \FunctionTok{head}\NormalTok{(}\DecValTok{15}\NormalTok{)}
\end{Highlighting}
\end{Shaded}

\begin{verbatim}
## ManufacturingProcess03 ManufacturingProcess11 ManufacturingProcess10 
##                     15                     10                      9 
## ManufacturingProcess25 ManufacturingProcess26 ManufacturingProcess27 
##                      5                      5                      5 
## ManufacturingProcess28 ManufacturingProcess29 ManufacturingProcess30 
##                      5                      5                      5 
## ManufacturingProcess31 ManufacturingProcess33 ManufacturingProcess34 
##                      5                      5                      5 
## ManufacturingProcess35 ManufacturingProcess36 ManufacturingProcess02 
##                      5                      5                      3
\end{verbatim}

\begin{Shaded}
\begin{Highlighting}[]
\NormalTok{chem\_pp }\OtherTok{\textless{}{-}} \FunctionTok{preProcess}\NormalTok{(x\_chem, }\AttributeTok{method =} \FunctionTok{c}\NormalTok{(}\StringTok{"knnImpute"}\NormalTok{, }\StringTok{"center"}\NormalTok{, }\StringTok{"scale"}\NormalTok{))}
\NormalTok{x\_chem\_imp }\OtherTok{\textless{}{-}} \FunctionTok{predict}\NormalTok{(chem\_pp, x\_chem)}
\FunctionTok{summary}\NormalTok{(}\FunctionTok{is.na}\NormalTok{(x\_chem\_imp))}
\end{Highlighting}
\end{Shaded}

\begin{verbatim}
##  BiologicalMaterial01 BiologicalMaterial02 BiologicalMaterial03
##  Mode :logical        Mode :logical        Mode :logical       
##  FALSE:176            FALSE:176            FALSE:176           
##  BiologicalMaterial04 BiologicalMaterial05 BiologicalMaterial06
##  Mode :logical        Mode :logical        Mode :logical       
##  FALSE:176            FALSE:176            FALSE:176           
##  BiologicalMaterial07 BiologicalMaterial08 BiologicalMaterial09
##  Mode :logical        Mode :logical        Mode :logical       
##  FALSE:176            FALSE:176            FALSE:176           
##  BiologicalMaterial10 BiologicalMaterial11 BiologicalMaterial12
##  Mode :logical        Mode :logical        Mode :logical       
##  FALSE:176            FALSE:176            FALSE:176           
##  ManufacturingProcess01 ManufacturingProcess02 ManufacturingProcess03
##  Mode :logical          Mode :logical          Mode :logical         
##  FALSE:176              FALSE:176              FALSE:176             
##  ManufacturingProcess04 ManufacturingProcess05 ManufacturingProcess06
##  Mode :logical          Mode :logical          Mode :logical         
##  FALSE:176              FALSE:176              FALSE:176             
##  ManufacturingProcess07 ManufacturingProcess08 ManufacturingProcess09
##  Mode :logical          Mode :logical          Mode :logical         
##  FALSE:176              FALSE:176              FALSE:176             
##  ManufacturingProcess10 ManufacturingProcess11 ManufacturingProcess12
##  Mode :logical          Mode :logical          Mode :logical         
##  FALSE:176              FALSE:176              FALSE:176             
##  ManufacturingProcess13 ManufacturingProcess14 ManufacturingProcess15
##  Mode :logical          Mode :logical          Mode :logical         
##  FALSE:176              FALSE:176              FALSE:176             
##  ManufacturingProcess16 ManufacturingProcess17 ManufacturingProcess18
##  Mode :logical          Mode :logical          Mode :logical         
##  FALSE:176              FALSE:176              FALSE:176             
##  ManufacturingProcess19 ManufacturingProcess20 ManufacturingProcess21
##  Mode :logical          Mode :logical          Mode :logical         
##  FALSE:176              FALSE:176              FALSE:176             
##  ManufacturingProcess22 ManufacturingProcess23 ManufacturingProcess24
##  Mode :logical          Mode :logical          Mode :logical         
##  FALSE:176              FALSE:176              FALSE:176             
##  ManufacturingProcess25 ManufacturingProcess26 ManufacturingProcess27
##  Mode :logical          Mode :logical          Mode :logical         
##  FALSE:176              FALSE:176              FALSE:176             
##  ManufacturingProcess28 ManufacturingProcess29 ManufacturingProcess30
##  Mode :logical          Mode :logical          Mode :logical         
##  FALSE:176              FALSE:176              FALSE:176             
##  ManufacturingProcess31 ManufacturingProcess32 ManufacturingProcess33
##  Mode :logical          Mode :logical          Mode :logical         
##  FALSE:176              FALSE:176              FALSE:176             
##  ManufacturingProcess34 ManufacturingProcess35 ManufacturingProcess36
##  Mode :logical          Mode :logical          Mode :logical         
##  FALSE:176              FALSE:176              FALSE:176             
##  ManufacturingProcess37 ManufacturingProcess38 ManufacturingProcess39
##  Mode :logical          Mode :logical          Mode :logical         
##  FALSE:176              FALSE:176              FALSE:176             
##  ManufacturingProcess40 ManufacturingProcess41 ManufacturingProcess42
##  Mode :logical          Mode :logical          Mode :logical         
##  FALSE:176              FALSE:176              FALSE:176             
##  ManufacturingProcess43 ManufacturingProcess44 ManufacturingProcess45
##  Mode :logical          Mode :logical          Mode :logical         
##  FALSE:176              FALSE:176              FALSE:176
\end{verbatim}

\textbf{Answer:} missing values are filled using k-nearest-neighbor
imputation.

\hypertarget{c-split-preprocess-and-tune-a-model-from-chapter-6}{%
\subsection{(c) Split, preprocess, and tune a model from Chapter
6}\label{c-split-preprocess-and-tune-a-model-from-chapter-6}}

My model choice here is \textbf{elastic net} because: - there may be
correlated process variables, - the sample size is not large, -
coefficient shrinkage can improve stability, - we still retain a
coefficient-based model that supports interpretation.

I also fit PLS as a benchmark because it is often competitive on
correlated process data.

\begin{Shaded}
\begin{Highlighting}[]
\FunctionTok{set.seed}\NormalTok{(}\DecValTok{2026}\NormalTok{)}
\NormalTok{idx\_chem }\OtherTok{\textless{}{-}} \FunctionTok{createDataPartition}\NormalTok{(y\_chem, }\AttributeTok{p =} \FloatTok{0.80}\NormalTok{, }\AttributeTok{list =} \ConstantTok{FALSE}\NormalTok{)}

\NormalTok{x\_chem\_train }\OtherTok{\textless{}{-}}\NormalTok{ x\_chem[idx\_chem, ]}
\NormalTok{x\_chem\_test  }\OtherTok{\textless{}{-}}\NormalTok{ x\_chem[}\SpecialCharTok{{-}}\NormalTok{idx\_chem, ]}
\NormalTok{y\_chem\_train }\OtherTok{\textless{}{-}}\NormalTok{ y\_chem[idx\_chem]}
\NormalTok{y\_chem\_test  }\OtherTok{\textless{}{-}}\NormalTok{ y\_chem[}\SpecialCharTok{{-}}\NormalTok{idx\_chem]}

\NormalTok{chem\_ctrl }\OtherTok{\textless{}{-}} \FunctionTok{trainControl}\NormalTok{(}\AttributeTok{method =} \StringTok{"cv"}\NormalTok{, }\AttributeTok{number =} \DecValTok{5}\NormalTok{)}

\NormalTok{chem\_pls }\OtherTok{\textless{}{-}} \FunctionTok{train}\NormalTok{(}
  \AttributeTok{x =}\NormalTok{ x\_chem\_train,}
  \AttributeTok{y =}\NormalTok{ y\_chem\_train,}
  \AttributeTok{method =} \StringTok{"pls"}\NormalTok{,}
  \AttributeTok{tuneGrid =} \FunctionTok{expand.grid}\NormalTok{(}\AttributeTok{ncomp =} \DecValTok{1}\SpecialCharTok{:}\DecValTok{15}\NormalTok{),}
  \AttributeTok{preProcess =} \FunctionTok{c}\NormalTok{(}\StringTok{"knnImpute"}\NormalTok{, }\StringTok{"center"}\NormalTok{, }\StringTok{"scale"}\NormalTok{),}
  \AttributeTok{trControl =}\NormalTok{ chem\_ctrl,}
  \AttributeTok{metric =} \StringTok{"RMSE"}
\NormalTok{)}

\NormalTok{chem\_enet }\OtherTok{\textless{}{-}} \FunctionTok{train}\NormalTok{(}
  \AttributeTok{x =}\NormalTok{ x\_chem\_train,}
  \AttributeTok{y =}\NormalTok{ y\_chem\_train,}
  \AttributeTok{method =} \StringTok{"enet"}\NormalTok{,}
  \AttributeTok{tuneGrid =} \FunctionTok{expand.grid}\NormalTok{(}\AttributeTok{lambda =} \FunctionTok{c}\NormalTok{(}\FloatTok{0.00}\NormalTok{, }\FloatTok{0.01}\NormalTok{, }\FloatTok{0.05}\NormalTok{), }\AttributeTok{fraction =} \FunctionTok{seq}\NormalTok{(}\FloatTok{0.10}\NormalTok{, }\DecValTok{1}\NormalTok{, }\AttributeTok{length.out =} \DecValTok{10}\NormalTok{)),}
  \AttributeTok{preProcess =} \FunctionTok{c}\NormalTok{(}\StringTok{"knnImpute"}\NormalTok{, }\StringTok{"center"}\NormalTok{, }\StringTok{"scale"}\NormalTok{),}
  \AttributeTok{trControl =}\NormalTok{ chem\_ctrl,}
  \AttributeTok{metric =} \StringTok{"RMSE"}
\NormalTok{)}

\NormalTok{chem\_lm }\OtherTok{\textless{}{-}} \FunctionTok{train}\NormalTok{(}
  \AttributeTok{x =}\NormalTok{ x\_chem\_train,}
  \AttributeTok{y =}\NormalTok{ y\_chem\_train,}
  \AttributeTok{method =} \StringTok{"lm"}\NormalTok{,}
  \AttributeTok{preProcess =} \FunctionTok{c}\NormalTok{(}\StringTok{"knnImpute"}\NormalTok{, }\StringTok{"center"}\NormalTok{, }\StringTok{"scale"}\NormalTok{),}
  \AttributeTok{trControl =}\NormalTok{ chem\_ctrl,}
  \AttributeTok{metric =} \StringTok{"RMSE"}
\NormalTok{)}

\NormalTok{chem\_models }\OtherTok{\textless{}{-}} \FunctionTok{list}\NormalTok{(}\AttributeTok{LM =}\NormalTok{ chem\_lm, }\AttributeTok{PLS =}\NormalTok{ chem\_pls, }\AttributeTok{ElasticNet =}\NormalTok{ chem\_enet)}
\FunctionTok{summary}\NormalTok{(}\FunctionTok{resamples}\NormalTok{(chem\_models))}
\end{Highlighting}
\end{Shaded}

\begin{verbatim}
## 
## Call:
## summary.resamples(object = resamples(chem_models))
## 
## Models: LM, PLS, ElasticNet 
## Number of resamples: 5 
## 
## MAE 
##                 Min.   1st Qu.    Median     Mean   3rd Qu.     Max. NA's
## LM         1.0546174 1.1606477 1.1771644 2.223575 1.2503733 6.475074    0
## PLS        0.8025465 1.0370804 1.0465846 1.026813 1.0888404 1.159012    0
## ElasticNet 0.6721988 0.8791256 0.8863139 0.874652 0.8933251 1.042296    0
## 
## RMSE 
##                 Min.  1st Qu.   Median     Mean  3rd Qu.      Max. NA's
## LM         1.2413009 1.352700 1.457929 6.714996 1.912527 27.610523    0
## PLS        1.1856506 1.287629 1.320378 1.297017 1.329736  1.361689    0
## ElasticNet 0.8196107 1.088340 1.118194 1.091072 1.124350  1.304864    0
## 
## Rsquared 
##                    Min.   1st Qu.    Median      Mean   3rd Qu.      Max. NA's
## LM         0.0001786201 0.3786757 0.5319395 0.4226357 0.5762273 0.6261574    0
## PLS        0.4640524696 0.5304580 0.5505597 0.5445897 0.5600670 0.6178111    0
## ElasticNet 0.5876051187 0.6074860 0.6370570 0.6736118 0.7337756 0.8021351    0
\end{verbatim}

\begin{Shaded}
\begin{Highlighting}[]
\FunctionTok{summary}\NormalTok{(}\FunctionTok{diff}\NormalTok{(}\FunctionTok{resamples}\NormalTok{(chem\_models)))}
\end{Highlighting}
\end{Shaded}

\begin{verbatim}
## 
## Call:
## summary.diff.resamples(object = diff(resamples(chem_models)))
## 
## p-value adjustment: bonferroni 
## Upper diagonal: estimates of the difference
## Lower diagonal: p-value for H0: difference = 0
## 
## MAE 
##            LM     PLS    ElasticNet
## LM                1.1968 1.3489    
## PLS        0.9697        0.1522    
## ElasticNet 0.8159 0.4050           
## 
## RMSE 
##            LM     PLS    ElasticNet
## LM                5.4180 5.6239    
## PLS        1.0000        0.2059    
## ElasticNet 1.0000 0.2386           
## 
## Rsquared 
##            LM     PLS    ElasticNet
## LM                -0.122 -0.251    
## PLS        0.8892        -0.129    
## ElasticNet 0.2464 0.1947
\end{verbatim}

\textbf{Answer:} the optimal value of the performance metric is the
\textbf{smallest cross-validated RMSE} among the tuned models. For the
chosen elastic-net model, the optimal tuning parameters are the selected
\texttt{lambda} and \texttt{fraction}.

\hypertarget{d-predict-the-test-set-and-compare-with-resampling}{%
\subsection{(d) Predict the test set and compare with
resampling}\label{d-predict-the-test-set-and-compare-with-resampling}}

\begin{Shaded}
\begin{Highlighting}[]
\NormalTok{chem\_test\_perf }\OtherTok{\textless{}{-}} \FunctionTok{bind\_rows}\NormalTok{(}\FunctionTok{lapply}\NormalTok{(}\FunctionTok{names}\NormalTok{(chem\_models), }\ControlFlowTok{function}\NormalTok{(nm) \{}
\NormalTok{  pred }\OtherTok{\textless{}{-}} \FunctionTok{predict}\NormalTok{(chem\_models[[nm]], x\_chem\_test)}
  \FunctionTok{tibble}\NormalTok{(}\AttributeTok{model =}\NormalTok{ nm,}
         \AttributeTok{RMSE =}\NormalTok{ caret}\SpecialCharTok{::}\FunctionTok{RMSE}\NormalTok{(pred, y\_chem\_test),}
         \AttributeTok{Rsquared =}\NormalTok{ caret}\SpecialCharTok{::}\FunctionTok{R2}\NormalTok{(pred, y\_chem\_test))}
\NormalTok{\})) }\SpecialCharTok{\%\textgreater{}\%} \FunctionTok{arrange}\NormalTok{(RMSE)}

\NormalTok{chem\_test\_perf}
\end{Highlighting}
\end{Shaded}

\begin{verbatim}
## # A tibble: 3 x 3
##   model       RMSE Rsquared
##   <chr>      <dbl>    <dbl>
## 1 ElasticNet  1.46   0.363 
## 2 PLS         3.03   0.136 
## 3 LM         10.4    0.0363
\end{verbatim}

\textbf{Answer:} compare the best model's test RMSE with its resampled
RMSE: - if they are close, the model generalized well; - if the test
RMSE is much worse, the model was probably overfit.

\hypertarget{e-which-predictors-are-most-important-do-biological-or-process-variables-dominate}{%
\subsection{(e) Which predictors are most important? Do biological or
process variables
dominate?}\label{e-which-predictors-are-most-important-do-biological-or-process-variables-dominate}}

For PLS, variable importance can be estimated through VIP; for the
elastic net, coefficient magnitude after standardization is an
interpretable proxy.

\hypertarget{pls-vip}{%
\subsubsection{PLS VIP}\label{pls-vip}}

\begin{Shaded}
\begin{Highlighting}[]
\NormalTok{chem\_vip }\OtherTok{\textless{}{-}} \FunctionTok{vip\_pls}\NormalTok{(chem\_pls)}
\FunctionTok{head}\NormalTok{(chem\_vip, }\DecValTok{15}\NormalTok{)}
\end{Highlighting}
\end{Shaded}

\begin{verbatim}
## # A tibble: 15 x 2
##    variable                 VIP
##    <chr>                  <dbl>
##  1 ManufacturingProcess32  2.04
##  2 ManufacturingProcess13  1.87
##  3 ManufacturingProcess17  1.75
##  4 ManufacturingProcess36  1.74
##  5 ManufacturingProcess09  1.72
##  6 BiologicalMaterial06    1.69
##  7 BiologicalMaterial02    1.64
##  8 BiologicalMaterial03    1.56
##  9 BiologicalMaterial08    1.41
## 10 BiologicalMaterial12    1.41
## 11 BiologicalMaterial11    1.39
## 12 ManufacturingProcess33  1.36
## 13 ManufacturingProcess06  1.35
## 14 BiologicalMaterial04    1.33
## 15 BiologicalMaterial01    1.32
\end{verbatim}

\hypertarget{elastic-net-standardized-coefficients}{%
\subsubsection{Elastic-net standardized
coefficients}\label{elastic-net-standardized-coefficients}}

\begin{Shaded}
\begin{Highlighting}[]
\NormalTok{chem\_coef }\OtherTok{\textless{}{-}} \FunctionTok{predict}\NormalTok{(chem\_enet}\SpecialCharTok{$}\NormalTok{finalModel,}
                     \AttributeTok{s =}\NormalTok{ chem\_enet}\SpecialCharTok{$}\NormalTok{bestTune}\SpecialCharTok{$}\NormalTok{fraction,}
                     \AttributeTok{mode =} \StringTok{"fraction"}\NormalTok{,}
                     \AttributeTok{type =} \StringTok{"coefficients"}\NormalTok{)}

\NormalTok{coef\_tbl }\OtherTok{\textless{}{-}} \FunctionTok{tibble}\NormalTok{(}
  \AttributeTok{variable =} \FunctionTok{names}\NormalTok{(chem\_coef}\SpecialCharTok{$}\NormalTok{coefficients),}
  \AttributeTok{coefficient =} \FunctionTok{as.numeric}\NormalTok{(chem\_coef}\SpecialCharTok{$}\NormalTok{coefficients)}
\NormalTok{) }\SpecialCharTok{\%\textgreater{}\%}
  \FunctionTok{filter}\NormalTok{(variable }\SpecialCharTok{!=} \StringTok{"(Intercept)"}\NormalTok{, coefficient }\SpecialCharTok{!=} \DecValTok{0}\NormalTok{) }\SpecialCharTok{\%\textgreater{}\%}
  \FunctionTok{mutate}\NormalTok{(}\AttributeTok{abs\_coefficient =} \FunctionTok{abs}\NormalTok{(coefficient)) }\SpecialCharTok{\%\textgreater{}\%}
  \FunctionTok{arrange}\NormalTok{(}\FunctionTok{desc}\NormalTok{(abs\_coefficient))}

\FunctionTok{head}\NormalTok{(coef\_tbl, }\DecValTok{15}\NormalTok{)}
\end{Highlighting}
\end{Shaded}

\begin{verbatim}
## # A tibble: 15 x 3
##    variable               coefficient abs_coefficient
##    <chr>                        <dbl>           <dbl>
##  1 ManufacturingProcess32      0.760           0.760 
##  2 ManufacturingProcess09      0.496           0.496 
##  3 ManufacturingProcess13     -0.376           0.376 
##  4 BiologicalMaterial06        0.272           0.272 
##  5 ManufacturingProcess37     -0.271           0.271 
##  6 ManufacturingProcess45      0.188           0.188 
##  7 ManufacturingProcess34      0.165           0.165 
##  8 ManufacturingProcess28     -0.163           0.163 
##  9 ManufacturingProcess39      0.147           0.147 
## 10 BiologicalMaterial05        0.147           0.147 
## 11 ManufacturingProcess06      0.138           0.138 
## 12 ManufacturingProcess15      0.134           0.134 
## 13 ManufacturingProcess07     -0.130           0.130 
## 14 ManufacturingProcess36     -0.122           0.122 
## 15 ManufacturingProcess01      0.0908          0.0908
\end{verbatim}

To separate biological from process variables, create a lookup vector
based on the original column names. If the package uses the first 12
columns for the biological measurements and the remaining 45 for process
variables, the code below will tag them accordingly.

\begin{Shaded}
\begin{Highlighting}[]
\NormalTok{chem\_var\_types }\OtherTok{\textless{}{-}} \FunctionTok{tibble}\NormalTok{(}
  \AttributeTok{variable =} \FunctionTok{colnames}\NormalTok{(x\_chem),}
  \AttributeTok{type =} \FunctionTok{c}\NormalTok{(}\FunctionTok{rep}\NormalTok{(}\StringTok{"Biological"}\NormalTok{, }\DecValTok{12}\NormalTok{), }\FunctionTok{rep}\NormalTok{(}\StringTok{"Process"}\NormalTok{, }\FunctionTok{ncol}\NormalTok{(x\_chem) }\SpecialCharTok{{-}} \DecValTok{12}\NormalTok{))}
\NormalTok{)}

\NormalTok{chem\_importance }\OtherTok{\textless{}{-}} \FunctionTok{full\_join}\NormalTok{(}
\NormalTok{  chem\_var\_types,}
\NormalTok{  coef\_tbl }\SpecialCharTok{\%\textgreater{}\%} \FunctionTok{select}\NormalTok{(variable, abs\_coefficient),}
  \AttributeTok{by =} \StringTok{"variable"}
\NormalTok{) }\SpecialCharTok{\%\textgreater{}\%} \FunctionTok{arrange}\NormalTok{(}\FunctionTok{desc}\NormalTok{(abs\_coefficient))}

\FunctionTok{head}\NormalTok{(chem\_importance, }\DecValTok{20}\NormalTok{)}
\end{Highlighting}
\end{Shaded}

\begin{verbatim}
## # A tibble: 20 x 3
##    variable               type       abs_coefficient
##    <chr>                  <chr>                <dbl>
##  1 ManufacturingProcess32 Process             0.760 
##  2 ManufacturingProcess09 Process             0.496 
##  3 ManufacturingProcess13 Process             0.376 
##  4 BiologicalMaterial06   Biological          0.272 
##  5 ManufacturingProcess37 Process             0.271 
##  6 ManufacturingProcess45 Process             0.188 
##  7 ManufacturingProcess34 Process             0.165 
##  8 ManufacturingProcess28 Process             0.163 
##  9 ManufacturingProcess39 Process             0.147 
## 10 BiologicalMaterial05   Biological          0.147 
## 11 ManufacturingProcess06 Process             0.138 
## 12 ManufacturingProcess15 Process             0.134 
## 13 ManufacturingProcess07 Process             0.130 
## 14 ManufacturingProcess36 Process             0.122 
## 15 ManufacturingProcess01 Process             0.0908
## 16 ManufacturingProcess17 Process             0.0838
## 17 ManufacturingProcess04 Process             0.0612
## 18 BiologicalMaterial07   Biological          0.0500
## 19 ManufacturingProcess20 Process             0.0466
## 20 ManufacturingProcess43 Process             0.0385
\end{verbatim}

\begin{Shaded}
\begin{Highlighting}[]
\FunctionTok{count}\NormalTok{(chem\_importance }\SpecialCharTok{\%\textgreater{}\%} \FunctionTok{slice\_head}\NormalTok{(}\AttributeTok{n =} \DecValTok{10}\NormalTok{), type)}
\end{Highlighting}
\end{Shaded}

\begin{verbatim}
## # A tibble: 2 x 2
##   type           n
##   <chr>      <int>
## 1 Biological     2
## 2 Process        8
\end{verbatim}

\textbf{Answer:} whichever group appears more often among the top-ranked
variables dominates the importance list. From a business perspective,
process variables are especially valuable because they are actionable.

\hypertarget{f-explore-the-relationships-between-top-predictors-and-yield}{%
\subsection{(f) Explore the relationships between top predictors and
yield}\label{f-explore-the-relationships-between-top-predictors-and-yield}}

\begin{Shaded}
\begin{Highlighting}[]
\NormalTok{top\_vars }\OtherTok{\textless{}{-}}\NormalTok{ chem\_importance }\SpecialCharTok{\%\textgreater{}\%}
  \FunctionTok{filter}\NormalTok{(}\SpecialCharTok{!}\FunctionTok{is.na}\NormalTok{(abs\_coefficient)) }\SpecialCharTok{\%\textgreater{}\%}
  \FunctionTok{slice\_head}\NormalTok{(}\AttributeTok{n =} \DecValTok{6}\NormalTok{) }\SpecialCharTok{\%\textgreater{}\%}
  \FunctionTok{pull}\NormalTok{(variable)}

\NormalTok{top\_vars}
\end{Highlighting}
\end{Shaded}

\begin{verbatim}
## [1] "ManufacturingProcess32" "ManufacturingProcess09" "ManufacturingProcess13"
## [4] "BiologicalMaterial06"   "ManufacturingProcess37" "ManufacturingProcess45"
\end{verbatim}

\begin{Shaded}
\begin{Highlighting}[]
\NormalTok{plot\_list }\OtherTok{\textless{}{-}} \FunctionTok{lapply}\NormalTok{(top\_vars, }\ControlFlowTok{function}\NormalTok{(v) \{}
  \FunctionTok{ggplot}\NormalTok{(}\FunctionTok{data.frame}\NormalTok{(}\AttributeTok{x =}\NormalTok{ x\_chem\_imp[[v]], }\AttributeTok{y =}\NormalTok{ y\_chem), }\FunctionTok{aes}\NormalTok{(}\AttributeTok{x =}\NormalTok{ x, }\AttributeTok{y =}\NormalTok{ y)) }\SpecialCharTok{+}
    \FunctionTok{geom\_point}\NormalTok{(}\AttributeTok{alpha =} \FloatTok{0.6}\NormalTok{, }\AttributeTok{shape =} \DecValTok{16}\NormalTok{) }\SpecialCharTok{+}
    \FunctionTok{geom\_smooth}\NormalTok{(}\AttributeTok{method =} \StringTok{"loess"}\NormalTok{, }\AttributeTok{se =} \ConstantTok{FALSE}\NormalTok{, }\AttributeTok{linewidth =} \FloatTok{0.4}\NormalTok{, }\AttributeTok{color =} \StringTok{"gray25"}\NormalTok{) }\SpecialCharTok{+}
    \FunctionTok{labs}\NormalTok{(}\AttributeTok{title =}\NormalTok{ v, }\AttributeTok{x =}\NormalTok{ v, }\AttributeTok{y =} \StringTok{"Yield"}\NormalTok{)}
\NormalTok{\})}

\NormalTok{plot\_list[[}\DecValTok{1}\NormalTok{]]}
\end{Highlighting}
\end{Shaded}

\begin{center}\includegraphics{chapter6_files/figure-latex/chem-top-plots-1} \end{center}

\begin{Shaded}
\begin{Highlighting}[]
\CommentTok{\# Print the rest one{-}by{-}one when knitting interactively:}
\NormalTok{plot\_list[[}\DecValTok{2}\NormalTok{]]}
\NormalTok{plot\_list[[}\DecValTok{3}\NormalTok{]]}
\NormalTok{plot\_list[[}\DecValTok{4}\NormalTok{]]}
\NormalTok{plot\_list[[}\DecValTok{5}\NormalTok{]]}
\NormalTok{plot\_list[[}\DecValTok{6}\NormalTok{]]}
\end{Highlighting}
\end{Shaded}

\textbf{Answer:} these plots are useful in two distinct ways.

\begin{enumerate}
\def\labelenumi{\arabic{enumi}.}
\tightlist
\item
  \textbf{Raw-material screening:} if a biological variable is strongly
  associated with yield, operators can use it to evaluate incoming
  material before production.
\item
  \textbf{Process optimization:} if a process variable shows a favorable
  region, the manufacturing team can target that operating range in
  future runs.
\end{enumerate}

So the practical value of the model is not only prediction; it is also
prioritizing controllable levers for yield improvement.

\hypertarget{final-answers-in-plain-english}{%
\section{Final answers in plain
English}\label{final-answers-in-plain-english}}

\hypertarget{tecator}{%
\subsection{6.1 Tecator}\label{tecator}}

\begin{itemize}
\tightlist
\item
  The spectra are highly collinear, so PCA should show a much smaller
  effective dimension than 100.
\item
  PLS is usually the most natural model choice because it handles
  correlated predictors directly and uses the response.
\item
  I would choose the model with the smallest test RMSE, but if several
  are tied, I would use PLS.
\end{itemize}

\hypertarget{permeability-1}{%
\subsection{6.2 Permeability}\label{permeability-1}}

\begin{itemize}
\tightlist
\item
  Remove sparse near-zero-variance fingerprints first.
\item
  Tune PLS as the primary latent-variable model.
\item
  Compare against ridge/lasso/elastic net.
\item
  I would recommend the best model as a \textbf{screening aid}, not a
  full replacement for the assay, unless the test-set error is
  exceptionally small.
\end{itemize}

\hypertarget{chemical-manufacturing}{%
\subsection{6.3 Chemical manufacturing}\label{chemical-manufacturing}}

\begin{itemize}
\tightlist
\item
  Impute missing values, then center/scale.
\item
  Fit at least PLS and elastic net.
\item
  Compare resampled RMSE to test RMSE to check generalization.
\item
  Inspect importance and top predictor-yield relationships.
\item
  If process variables dominate the top of the list, that is especially
  useful because those variables are controllable.
\end{itemize}

\hypertarget{session-info}{%
\section{Session info}\label{session-info}}

\begin{Shaded}
\begin{Highlighting}[]
\FunctionTok{sessionInfo}\NormalTok{()}
\end{Highlighting}
\end{Shaded}

\begin{verbatim}
## R version 4.2.2 Patched (2022-11-10 r83330)
## Platform: x86_64-pc-linux-gnu (64-bit)
## Running under: Debian GNU/Linux 12 (bookworm)
## 
## Matrix products: default
## BLAS:   /usr/lib/x86_64-linux-gnu/blas/libblas.so.3.11.0
## LAPACK: /usr/lib/x86_64-linux-gnu/lapack/liblapack.so.3.11.0
## 
## locale:
##  [1] LC_CTYPE=C.UTF-8       LC_NUMERIC=C           LC_TIME=C.UTF-8       
##  [4] LC_COLLATE=C.UTF-8     LC_MONETARY=C.UTF-8    LC_MESSAGES=C.UTF-8   
##  [7] LC_PAPER=C.UTF-8       LC_NAME=C              LC_ADDRESS=C          
## [10] LC_TELEPHONE=C         LC_MEASUREMENT=C.UTF-8 LC_IDENTIFICATION=C   
## 
## attached base packages:
## [1] stats     graphics  grDevices utils     datasets  methods   base     
## 
## other attached packages:
##  [1] knitr_1.51                      tidyr_1.3.2                    
##  [3] dplyr_1.2.0                     pROC_1.19.0.1                  
##  [5] pls_2.9-0                       elasticnet_1.3                 
##  [7] lars_1.3                        AppliedPredictiveModeling_1.1-7
##  [9] caret_7.0-1                     lattice_0.20-45                
## [11] ggplot2_4.0.2                  
## 
## loaded via a namespace (and not attached):
##  [1] Rcpp_1.1.1           rpart.plot_3.1.4     lubridate_1.9.5     
##  [4] listenv_0.10.1       class_7.3-21         utf8_1.2.6          
##  [7] digest_0.6.39        ipred_0.9-15         foreach_1.5.2       
## [10] parallelly_1.46.1    R6_2.6.1             plyr_1.8.9          
## [13] hardhat_1.4.2        stats4_4.2.2         ellipse_0.5.0       
## [16] evaluate_1.0.5       pillar_1.11.1        rlang_1.1.7         
## [19] data.table_1.18.2.1  rpart_4.1.19         Matrix_1.5-3        
## [22] rmarkdown_2.31       labeling_0.4.3       splines_4.2.2       
## [25] gower_1.0.2          stringr_1.6.0        S7_0.2.1            
## [28] compiler_4.2.2       xfun_0.56            pkgconfig_2.0.3     
## [31] mgcv_1.8-41          globals_0.19.1       htmltools_0.5.9     
## [34] nnet_7.3-18          tidyselect_1.2.1     tibble_3.3.1        
## [37] prodlim_2026.03.11   RANN_2.6.2           CORElearn_1.57.3.1  
## [40] codetools_0.2-19     future_1.70.0        withr_3.0.2         
## [43] MASS_7.3-58.2        recipes_1.3.1        ModelMetrics_1.2.2.2
## [46] grid_4.2.2           nlme_3.1-162         gtable_0.3.6        
## [49] lifecycle_1.0.5      magrittr_2.0.4       scales_1.4.0        
## [52] future.apply_1.20.2  cli_3.6.5            stringi_1.8.7       
## [55] farver_2.1.2         reshape2_1.4.5       timeDate_4052.112   
## [58] generics_0.1.4       vctrs_0.7.1          lava_1.8.2          
## [61] RColorBrewer_1.1-3   iterators_1.0.14     tools_4.2.2         
## [64] glue_1.8.0           purrr_1.2.1          plotrix_3.8-14      
## [67] parallel_4.2.2       fastmap_1.2.0        survival_3.5-3      
## [70] yaml_2.3.12          timechange_0.4.0     cluster_2.1.4
\end{verbatim}

\end{document}
